\chapter{不变卡尔曼滤波与流形滤波族}
\label{ch:invariant-filter}

% ════════════════════════════════════════════════════════════════
%  本章吸收：Tutorial 60_概率与估计/30_流形滤波族.md（ESKF/MEKF/IEKF/InEKF/UKF-M 谱系，
%  Solà-ESKF、Lefferts-Markley-Shuster、Hartley contact-InEKF、Brossard UKF-M）
%  + 70_Barrau_Bonnabel精读.md（TAC2017 群仿射 Thm1、log-linear Thm2、稳定性 Thm4 逐定理证明）。
%  承接 ch:eskf（ESKF 主线 + paper:inekf 概述）与 ch:lie（SE_2(3) 权威定义 sec:se23、
%  对称/不变/等变 sec:lie-symmetry）。右扰动为主；SE_2(3) 群元/漂移矩阵按本书 def:se23 [rho;ν;φ]，
%  切空间系统/观测矩阵随源文献 [deltaθ;deltav;deltap]（旋转在前），note 盒显式对照置换。教学层全要素（v5.0）。
% ════════════════════════════════════════════════════════════════

\begin{practice}[前置自测]
开始之前，先试答下列问题。若能流畅作答，可只速读本章表格与速查卡；若卡住，正是本章要为你解决的。
\begin{enumerate}
  \item 把四元数 $\mathbf{q}\in\mathbb{H}$ 的 $4$ 个分量直接塞进卡尔曼滤波（KF）的状态向量、维护 $4\times4$ 协方差，会出什么问题？为什么？
  \item \cref{ch:eskf} 的误差状态卡尔曼滤波（ESKF）把状态拆成“标称 $+$ 误差 $+$ 真值”三态，是为了\emph{省计算}吗？根本动机是什么？
  \item 标准扩展卡尔曼滤波（EKF）的传播雅可比 $\mathbf{A}(\hat{\mathbf{x}},\mathbf{u})$ 依赖\emph{估计}状态 $\hat{\mathbf{x}}$。这会在 VIO/SLAM 里诱发什么后果（提示：yaw 漂移）？
  \item 给定李群上的两条轨迹 $\mathbf{X}_t,\hat{\mathbf{X}}_t$，“右不变误差” $\boldsymbol{\eta}=\hat{\mathbf{X}}\mathbf{X}^{-1}$ 与标准 EKF 的加性误差 $\mathbf{e}=\hat{\mathbf{x}}-\mathbf{x}$ 有何本质区别？
  \item 同样是 “IEKF”，有人指 \emph{iterated} EKF、有人指 \emph{invariant} EKF。你能分清这两者各自解决什么问题吗？
\end{enumerate}
\end{practice}

\section*{本章目标}
学完本章，你应当能够：
\begin{enumerate}
  \item \textbf{说清}把旋转/位姿当欧氏向量塞进 EKF 的三类\emph{结构性}失败（过参数化协方差奇异、参数化奇点、虚假可观测性），并据此理解“流形滤波”的必要性；
  \item \textbf{把 MEKF、ESKF、InEKF、UKF-M 统一}为同一原则——“让误差活在切空间、状态活在流形、把大运动交给标称态、把不确定度交给切空间扰动”——的四种局部坐标实现；
  \item \textbf{独立复现}从乘性姿态误差 $\mathbf{R}=\hat{\mathbf{R}}\,\mathrm{Exp}(\delta\boldsymbol{\theta})$ 推出 MEKF/ESKF 角误差与速度误差动力学的逐步推导；
  \item \textbf{陈述并证明} Barrau--Bonnabel 群仿射（group-affine）条件与左/右不变误差自治的三等价定理（\cref{ch:lie} 已给概念地基，本章给完整代数）；
  \item \textbf{解释} log-linear 定理为何是\emph{精确}而非近似——指出 BCH 交叉项如何经“流的群同态 $+\,[\mathbf{X},\mathbf{X}]=\mathbf{0}$”恰好相消；
  \item \textbf{逐块验证} $\mathrm{SE}_2(3)$ 上 IMU 动力学的群仿射性，并写出仅依赖输入的误差转移矩阵 $\mathbf{A}_t$（严格幂零，$\mathbf{A}_t^3=\mathbf{0}$）；
  \item \textbf{实现}右不变 InEKF（RIEKF）的完整离散预测—更新循环，正确处理“不变创新量、群乘法更新、左右手性、Joseph 形式协方差”；
  \item \textbf{论证} InEKF 为何\emph{结构性}地保持 yaw/平移不可观子空间（$\mathbf{H}\mathbf{N}=\mathbf{0}$ 在传播下保持），从而无需 FEJ/OC-EKF 的经验修补；
  \item \textbf{说明} IMU 零偏为何\emph{破坏}群仿射（negative result 的代数证明），以及“不完美 InEKF / 切群 / 等变滤波 EqF”三条修复路径；
  \item \textbf{掌握} UKF-M 的免雅可比 sigma 点流形滤波，区分它与 InEKF 的互补关系；分清 IKFoM（iterated）与 InEKF（invariant）的命名冲突；
  \item \textbf{陈述} InEKF 渐近稳定性定理（收敛半径 $\varepsilon$ 独立于初始化时刻 $t_0$）的条件与意义，理解它相对标准 EKF 稳定性结论的根本进步。
\end{enumerate}

\subsection*{本章知识导航}
本章是一条“\emph{把 KF 从欧氏空间搬到流形、再用对称性把误差动力学与状态解耦}”的主线。结构如下：

\begin{center}
\small
\begin{tikzpicture}[
  node distance=4.5mm and 7mm, >=Stealth,
  blk/.style={draw, rounded corners, align=center, font=\small, inner sep=3pt, fill=main!6, draw=main!60!black},
  grp/.style={draw, rounded corners, align=center, font=\small, inner sep=3pt, fill=second!6, draw=second!60!black},
  thy/.style={draw, rounded corners, align=center, font=\small, inner sep=3pt, fill=third!8, draw=third!60!black},
  ln/.style={->, thick, gray!60!black}]
\node[blk] (fail) {三类失败\\(奇异/奇点/伪观测)};
\node[blk, right=of fail] (split) {名义$+$误差\\切空间高斯};
\node[grp, right=of split] (mekf) {MEKF / ESKF\\乘性姿态误差};
\node[thy, below=8mm of fail] (ga) {群仿射\\$f(ab)$ 恒等式};
\node[thy, right=of ga] (loglin) {log-linear\\$\dot{\boldsymbol{\xi}}=\mathbf{A}_t\boldsymbol{\xi}$ 精确};
\node[thy, right=of loglin] (inekf) {InEKF\\$\mathrm{SE}_2(3)$ 不变误差};
\node[grp, below=8mm of ga] (stab) {稳定性\\$\varepsilon\perp t_0$};
\node[grp, right=of stab] (obs) {一致性\\$\mathbf{H}\mathbf{N}=\mathbf{0}$ 保持};
\node[blk, right=of obs, fill=third!8, draw=third!60!black] (ext) {UKF-M / IKFoM\\EqF（推广）};
\draw[ln] (fail)--(split); \draw[ln] (split)--(mekf);
\draw[ln] (mekf.south) -- ++(0,-3.5mm) -| (inekf.north);
\draw[ln] (ga)--(loglin); \draw[ln] (loglin)--(inekf);
\draw[ln] (ga.south)-- (stab.north);
\draw[ln] (stab)--(obs); \draw[ln] (obs)--(ext);
\end{tikzpicture}
\end{center}

\noindent\textbf{推荐路径}：\cref{sec:why-manifold,sec:family-unified} 是任何读者的主干（建立“为什么要流形滤波 $+$ 统一原则”）；\cref{sec:mekf-eskf} 把 \cref{ch:eskf} 的 ESKF 接到 MEKF 历史原型；\cref{sec:group-affine,sec:log-linear} 是\emph{理论核心}（群仿射与 log-linear，承接 \cref{sec:lie-symmetry} 的对称性地基），其证明在 \cref{ch:eskf} 的 \cref{paper:inekf} 只给结论、这里给完整代数；\cref{sec:inekf-algo,sec:consistency-inekf} 落到 InEKF 算法与一致性；\cref{sec:bias-negative,sec:ukfm-ikfom,sec:stability} 处理零偏边界、免雅可比变体与稳定性定理。带（选读）的进阶节可在首读时跳过。

\subsection*{前置知识桥接}
本章\emph{重度}依赖前面三章，请确认已掌握：
\begin{itemize}
  \item \cref{ch:lie}：李群 $\SO(3)$、$\SE(3)$、扩展位姿群 $\mathrm{SE}_2(3)$（权威定义 \cref{def:se23}、指数映射 \cref{thm:se23-exp}、伴随 \cref{eq:se23-adjoint}），指数/对数映射、左右雅可比、伴随 $\mathrm{Ad}$，以及\emph{对称、不变、等变}的群论语言（\cref{sec:lie-symmetry}、不变误差 \cref{def:invariant-error}）——这些是 InEKF 的几何与代数地基；
  \item \cref{ch:eskf}：卡尔曼滤波（\cref{thm:kf}）、扩展卡尔曼滤波（\cref{prop:ekf}）、误差状态滤波 ESKF（\cref{sec:eskf}、误差动力学 \cref{thm:eskf-error-kin}、复位 \cref{eq:eskf-reset}）、滤波器谱系（\cref{tab:filter-taxo}）、一致性诊断 NEES/NIS（\cref{def:nees-nis}）、以及不变滤波概览盒（\cref{paper:inekf}）；
  \item \cref{ch:slam_est}：MAP/最小二乘、信息矩阵、可观测性的概率视角。
\end{itemize}
\cref{ch:eskf} 的 \cref{paper:inekf} 已给出 InEKF 的“一段话概览”；\textbf{本章是该主题的完整、自包含展开}——群仿射三等价的逐步证明、log-linear 的 BCH 机理、$\mathrm{SE}_2(3)$ 逐块验证、离散算法、零偏 negative result、稳定性定理，以及整个流形滤波族（MEKF/ESKF/InEKF/UKF-M/IKFoM/EqF）的横向对比。

\subsection*{如果跳过本章会怎样}
\begin{itemize}
  \item 你会知道 ESKF（\cref{sec:eskf}）改善了一致性，却\emph{不知道}为什么它仍未根除“线性化误差与状态相关”这一病根，也无法判断何时该升级到 InEKF；
  \item 在 VIO/腿足里程计里，你会照抄某个开源 InEKF 库，却分不清左/右不变误差、左/右乘更新、$\mathrm{SE}_2(3)$ 切空间排序——这正是流形滤波\emph{集成 bug 的头号来源}；
  \item 读论文时，你会被 “IEKF” 一词在 \emph{iterated} 与 \emph{invariant} 之间反复绊倒，无法理解 FAST-LIO2（IKFoM）与 InEKF 在理论上的根本区别。
\end{itemize}

\begin{table}[H]\centering\small
\caption{本章阅读方式建议}
\begin{tabular}{lll}
\toprule
阅读方式 & 时间 & 适合谁 \\
\midrule
精读（含群仿射/log-linear/稳定性全部证明与练习） & 6--8 小时 & 要做 VIO/惯导/腿足一致性研究 \\
速读（跳过 \cref{sec:log-linear,sec:stability} 的证明细节） & 2.5 小时 & 有 ESKF 基础、想掌握 InEKF 算法与边界 \\
速查（只看小结的符号表 $+$ 定理速查表 $+$ 全景对比） & 20 分钟 & 写代码、选型时回来查 \\
\bottomrule
\end{tabular}
\end{table}

\begin{note}
\textbf{本章记号与约定.}\;
状态用李群元素 $\mathbf{X}\in G$（源文献\cite{barrau2017invariant} 常记 $\chi$，本书统一 $\mathbf{X}$），旋转 $\mathbf{R}\in\SO(3)$（源用 $\mathbf{C}$/$R$，本书统一 $\mathbf{R}$），位姿 $\mathbf{T}\in\SE(3)$，扩展位姿 $\mathbf{X}\in\mathrm{SE}_2(3)$。
$\mathrm{Exp}(\cdot):=\exp((\cdot)^\wedge)$、$\mathrm{Log}(\cdot):=\ln(\cdot)^\vee$（沿用 \cref{ch:lie}）。
\textbf{扰动方向：全书右扰动 $\mathbf{X}\,\mathrm{Exp}(\delta\boldsymbol{\xi})$ 为主}。但 InEKF 的\emph{核心}恰是左/右\emph{不变误差}并用（见 \cref{sec:group-affine}）——这是该方法的内禀结构、非约定选择；本章对二者均给出，并在更新步明确手性。
\textbf{四元数 Hamilton 约定}（与 Eigen/十四讲一致）；引用 JPL 来源（如 MEKF 航天文献、OpenVINS）处显式注明。
\textbf{$\mathrm{SE}_2(3)$ 分量排序（务必看清）.}\;全书权威定义 \cref{def:se23} 取 $\boldsymbol{\xi}=[\boldsymbol{\rho};\boldsymbol{\nu};\boldsymbol{\phi}]$（位置、速度、\textbf{旋转在最后}），对应伴随/系统矩阵分块\emph{上}三角；Barrau--Bonnabel/Hartley 源文献用 $[\boldsymbol{\phi};\boldsymbol{\nu}_v;\boldsymbol{\nu}_p]$（\textbf{旋转在最前}），分块\emph{下}三角。二者仅差一个置换 $\mathbf{P}$、是同一线性映射。\textbf{本章采用两层对照写法}：\emph{群元素与漂移矩阵} $f(\mathbf{X})$（\cref{eq:imu-matrix-f}）按本书 \cref{def:se23} 的列序书写（位置列在前，见该处说明）；而\emph{切空间的系统矩阵} $\mathbf{A}_t$\emph{、状态转移} $\boldsymbol{\Phi}$\emph{、观测雅可比} $\mathbf{H}$ 统一采用\textbf{源文献序} $\boldsymbol{\xi}=[\delta\boldsymbol{\theta};\delta\mathbf{v};\delta\mathbf{p}]$（旋转在前、\emph{下}三角），以便与 \cite{barrau2017invariant} 的原始矩阵\emph{逐块对照}、降低集成 bug；需与 $[\boldsymbol{\rho};\boldsymbol{\nu};\boldsymbol{\phi}]$ 序对接时按置换 $\mathbf{P}$ 转换即可。
\end{note}

% ════════════════════════════════════════════════════════════════
\section{为什么需要流形滤波：三类结构性失败}
\label{sec:why-manifold}

\paragraph{动机：旋转不是向量，硬塞进 KF 会“结构性”崩。}
\cref{ch:eskf} 已建立卡尔曼滤波的完整机器，并在 \cref{sec:eskf} 用 ESKF 处理了带姿态的状态。本章把视野扩大到\emph{整个}流形滤波族，先回答一个更根本的问题：当状态含旋转/位姿时，\emph{为什么}不能像处理 $\mathbb{R}^n$ 那样直接套用 EKF？答案不是“精度差一点”，而是三类\emph{结构性}（非数值）失败\cite{sola2017quaternion}。

\paragraph{反面失败一：过参数化 $\Rightarrow$ 协方差奇异。}
四元数 $\mathbf{q}\in\mathbb{H}$ 有 $4$ 个分量，但单位范数约束 $\mathbf{q}^\top\mathbf{q}=1$ 使其切空间只有 $3$ 维。若用 $4\times4$ 协方差描述姿态不确定度，沿 $\mathbf{q}$ 径向（违反范数约束的方向）的协方差\emph{必然为零}——矩阵奇异，Cholesky 分解崩溃。这正是 Lefferts--Markley--Shuster 在 $1982$ 年提出乘性 EKF（MEKF）的直接动因\cite{lefferts1982kalman}。

\paragraph{反面失败二：参数化奇点 $\Rightarrow$ 雅可比病态。}
欧拉角在万向锁（gimbal lock）、轴角在 $\theta\to\pi$、修正罗德里格斯参数（MRP）在 $2\pi$ 附近都有不可微点。滤波器一旦穿越奇点，雅可比病态、协方差爆炸。这是 \cref{ch:eskf} 的 \cref{sec:eskf} 用最小 $3$ 参误差 $\delta\boldsymbol{\theta}$（而非全局参数）滤波的理由。

\paragraph{反面失败三：不可观方向“幽灵观测” $\Rightarrow$ 一致性退化。}
重力对齐的 VIO 沿典型轨迹有 $4$ 维不可观子空间（绝对位置 $3$ 维 $+$ 绕重力轴 yaw $1$ 维）；纯 $3$D SLAM 的全局规范（gauge）是 $6$ 维刚体变换。若雅可比在估计 $\hat{\mathbf{X}}$（而非真值）处求值，线性化后的可观测性 Gramian 可能\emph{丢掉} yaw 方向——滤波器“错觉”自己观测到了 yaw，协方差过度自信，长期 yaw 漂移\cite{sola2017quaternion}。这是 EKF-SLAM/VIO 二十年的一致性顽疾，催生了 Huang--Mourikis--Roumeliotis 的首次估计雅可比（FEJ）与可观测性约束 EKF（OC-EKF）等\emph{经验}修补，也是本章 InEKF 要从\emph{第一性原理}根治的对象。

\paragraph{本章导引。}
MEKF/ESKF（\cref{sec:mekf-eskf}）解决前两类问题；InEKF（\cref{sec:group-affine,sec:inekf-algo}）从对称性消除第三类问题；UKF-M/IKFoM/EqF（\cref{sec:ukfm-ikfom}）把这些思想扩展到\emph{免雅可比}与\emph{齐性空间}。

\begin{figure}[ht]
\centering
\begin{tikzpicture}[scale=1.0,>=Stealth,
  ax/.style={->, thick, gray!55!black},
  pt/.style={circle, fill=second, inner sep=1.3pt},
  ev/.style={font=\scriptsize}]
  % manifold (a wavy curve) = nominal trajectory
  \draw[thick, main] (-0.3,0) .. controls (1.2,1.0) and (2.8,-0.6) .. (4.3,0.5)
      node[right, main, font=\small] {流形 $\mathcal{M}$（标称轨迹）};
  % a base point with tangent line + gaussian
  \coordinate (P) at (1.85,0.18);
  \fill[main] (P) circle (1.2pt);
  \node[ev, below right=0pt and 1pt of P] {$\hat{\mathbf{X}}_k$};
  \draw[second, thick] (0.9,-0.55) -- (2.8,0.9) node[right, second, font=\scriptsize] {切空间 $T_{\hat{\mathbf{X}}_k}\mathcal{M}$};
  % gaussian blob on the tangent line
  \begin{scope}[shift={(P)}, rotate=37]
    \draw[third, thick, fill=third!12, fill opacity=0.5] (0,0) ellipse (0.62 and 0.22);
    \node[third, font=\scriptsize] at (0.0,0.42) {$\boldsymbol{\xi}\sim\mathcal{N}(\mathbf 0,\mathbf P)$};
  \end{scope}
  % exp arrow back to manifold
  \draw[->, third, thick] (2.55,0.78) .. controls (2.7,0.45) .. (2.62,0.15);
  \node[third, ev] at (3.15,0.55) {$\mathrm{Exp}$（注入）};
\end{tikzpicture}
\caption{流形滤波的统一图像：标称态 $\hat{\mathbf{X}}_k$ 是\emph{流形上}的一点、承载大运动；不确定度 $\boldsymbol{\xi}\sim\mathcal{N}(\mathbf 0,\mathbf P)$ 是其\emph{切空间}（欧氏空间）里的高斯，KF 在切空间做线性运算；每步更新后用指数映射把校正“注入”回流形。MEKF/ESKF/InEKF/UKF-M 是这同一图像的四种局部坐标选择。（仿 \cite{sola2017quaternion} 重绘）}
\label{fig:manifold-filter}
\end{figure}

\begin{pitfall}[概念误区：“归一化一下不就行了？”]
新手常想：把四元数当 $\mathbb{R}^4$ 做 EKF，每步\emph{归一化}回单位球面，不就解决了？归一化只修复了\emph{模值}约束，但 $4\times4$ 协方差沿径向仍有零特征值，Cholesky 必然奇异；更隐蔽的是，协方差\emph{并不知道}归一化改变了状态坐标——它仍以为自己活在 $4$ 维，而真实姿态只有 $3$ 维自由度，长期运行会在约束方向产生不可解释的数值退化。正解是直接用 $3$ 维误差协方差（MEKF/ESKF），而非在 $4$ 维空间打补丁（\cref{ch:eskf} 的 \cref{sec:eskf} 已点出此原则）。
\end{pitfall}

\begin{practice}
\begin{enumerate}
  \item 写出四元数单位范数约束 $\mathbf{q}^\top\mathbf{q}=1$，对其求导说明姿态扰动只能有 $3$ 个自由度，从而 $4\times4$ 协方差必沿某方向退化。
  \item 举例说明欧拉角万向锁：当俯仰角 $=90^\circ$ 时，哪两个轴的旋转变得不可区分？
  \item （概念）用一句话概括三类失败各自由哪类方法解决（对照本节“本章导引”）。
\end{enumerate}
\end{practice}

% ════════════════════════════════════════════════════════════════
\section{统一原则：同一件事的四种局部坐标}
\label{sec:family-unified}

\paragraph{动机：别把 MEKF/ESKF/InEKF/UKF-M 当四个互不相干的滤波器。}
初学者容易把这四种方法当成四套独立算法。更准确的理解是：它们都在回答\emph{同一个}问题。

\begin{insight}[流形滤波的一句话本质]
滤波器真正估计的不是“状态本身”，而是“\emph{真状态相对当前参考状态的小扰动}”。状态在流形上时，小扰动必须活在切空间；\emph{参考状态（标称态）承载大运动，切空间扰动承载不确定性}。这一观点与线性 KF 的区别极小，却足以改变整个工程实现。
\end{insight}

在线性 KF 中状态空间是 $\mathbb{R}^n$，真值、估计、误差自然满足 $\mathbf{x}=\hat{\mathbf{x}}+\delta\mathbf{x}$。这条式子背后有个隐含假设：任意两状态可\emph{相减}、差仍在同一线性空间。位置、速度满足；旋转不满足——两个旋转矩阵 $\mathbf{R}_1,\mathbf{R}_2$ 的差 $\mathbf{R}_1-\mathbf{R}_2$ 一般不是旋转，甚至无清晰物理意义（\cref{ch:lie} 已强调）。于是在李群上要把“减法”换成“相对变换 $+$ 取对数坐标”：
\begin{equation}
  \delta\mathbf{R}=\hat{\mathbf{R}}^\top\mathbf{R}\in\SO(3),\qquad
  \delta\boldsymbol{\theta}=\mathrm{Log}(\hat{\mathbf{R}}^\top\mathbf{R})\in\mathbb{R}^3.
  \label{eq:boxminus-rot}
\end{equation}
这就是 $\boxminus$ 算子的本质：不是普通减法，而是“先求相对位姿、再取对数坐标”。对应的“加法”$\boxplus$ 是 $\mathbf{R}\boxplus\delta\boldsymbol{\theta}=\mathbf{R}\,\mathrm{Exp}(\delta\boldsymbol{\theta})$（右扰动，\cref{ch:lie} 的现代语言）。

\begin{table}[H]\centering\small
\caption{四种流形滤波器：同一原则的四种局部坐标实现}
\label{tab:family-unified}
\begin{tabular}{lllll}
\toprule
方法 & 大状态在哪里传播 & 小误差如何定义 & KF 操作对象 & 关键收益 \\
\midrule
MEKF & $\SO(3)$ 姿态（四元数） & $\mathbf{q}=\hat{\mathbf{q}}\otimes\delta\mathbf{q}$（乘性） & $\delta\boldsymbol{\theta},\mathbf{b}_g$ & 姿态最小维度 \\
ESKF & $\SE(3)$/IMU 标称态 & 姿态乘性、其余加性 & $\delta\mathbf{p},\delta\mathbf{v},\delta\boldsymbol{\theta},\delta\mathbf{b}$ & IMU 工程主流 \\
InEKF & 矩阵李群 $G$ & $\boldsymbol{\eta}=\hat{\mathbf{X}}\mathbf{X}^{-1}$ 或 $\mathbf{X}^{-1}\hat{\mathbf{X}}$ & $\boldsymbol{\xi}=\mathrm{Log}(\boldsymbol{\eta})$ & 误差动力学与状态解耦 \\
UKF-M & 可平行化流形 $\mathcal{M}$ & retraction $\varphi,\varphi^{-1}$ & sigma 点切空间坐标 & 免解析雅可比 \\
\bottomrule
\end{tabular}
\end{table}

\paragraph{反面：不做这种拆分会怎样。}
标准 EKF 把旋转当欧氏向量、预测后归一化四元数。短时间内能跑，但协方差\emph{不知道}归一化改变了状态坐标，仍以为自己活在 $4$ 维，长期在约束方向数值退化（\cref{sec:why-manifold} 的失败一）。工程上这与机器人\emph{局部地图}维护异曲同工：全局地图很大，但 scan-to-map 每次只在当前位姿附近求一个小增量——大地图承载全局结构、小增量承载快速优化；ESKF/MEKF/InEKF 同理，\emph{标称态承载大运动、误差态承载小校正}。

\begin{note}
\textbf{与 \cref{ch:eskf} 的关系.}\;
\cref{ch:eskf} 的 \cref{sec:eskf} 已把 ESKF 讲透（三态、误差动力学 \cref{thm:eskf-error-kin}、复位 \cref{eq:eskf-reset}），\cref{tab:filter-taxo} 已给滤波器谱系。本章不重复 ESKF 全部细节，而是：(i) 把 ESKF 接到它的历史原型 MEKF（\cref{sec:mekf-eskf}）；(ii) 把 \cref{paper:inekf} 的 InEKF 概览展开为完整理论与算法（\cref{sec:group-affine,sec:inekf-algo}）；(iii) 补齐 UKF-M/IKFoM/EqF 与全族对比。读者可把本章视为 \cref{ch:eskf} 在“流形 $+$ 对称性”维度上的纵深。
\end{note}

\begin{practice}
\begin{enumerate}
  \item 在平移群 $\mathbb{R}$ 上，$\boxminus$ 退化为普通减法吗？在 $\SO(2)$ 上呢（提示：考虑 $170^\circ$ 与 $-170^\circ$ 的“差”应是 $20^\circ$ 还是 $340^\circ$）？
  \item 对照 \cref{tab:family-unified}，指出 ESKF 与 InEKF 在“误差定义”上的关键差异（加性 vs 群乘性），并预测哪个的误差动力学更可能与状态解耦。
\end{enumerate}
\end{practice}

% ════════════════════════════════════════════════════════════════
\section{MEKF 与 ESKF：乘性姿态误差的历史与推导}
\label{sec:mekf-eskf}

\paragraph{动机：ESKF 不是凭空出现，它的祖先是 1982 年的 MEKF。}
\cref{ch:eskf} 的 \cref{sec:eskf} 直接给出了 ESKF 的 $18$ 维误差动力学。本节补上它的\emph{历史脉络}与\emph{最小推导}：乘性 EKF（MEKF）是把“乘性姿态误差”引入 KF 的开山之作，ESKF 不过是它在“IMU 驱动的 $\SE(3)$ $+$ 双零偏 $+$ 重力”上的自然推广。

\subsection{MEKF：\texorpdfstring{$\SO(3)$}{SO(3)} 上的乘性误差原型}
\label{subsec:mekf}

\paragraph{核心思想（Lefferts--Markley--Shuster 1982）。}
把真姿态写成参考四元数（标称）乘上小角误差\cite{lefferts1982kalman}：$\mathbf{q}=\bar{\mathbf{q}}\otimes\delta\mathbf{q}(\delta\boldsymbol{\alpha})$，KF 只操作 $3$ 维小角误差 $\delta\boldsymbol{\alpha}$ 与陀螺零偏 $\mathbf{b}_g$（共 $6$ 维），彻底避开四元数过参数化奇异（\cref{sec:why-manifold} 失败一）。航天领域标准选 JPL 约定（global-to-local，左手代数 $ji=k$）。

\begin{note}
\textbf{记号对照.}\;
本书用 Hamilton 约定（$ij=k$，与 Eigen/十四讲一致）；MEKF 原始文献\cite{lefferts1982kalman} 及多数航天 AOCS 代码用 JPL 约定（$ji=k$）。同一组 $(x,y,z,w)$ 数值在两约定下\emph{语义相反}（OpenVINS 官方文档亦专门提醒）；跨库对接必须加显式转换层。本节推导用 Hamilton。
\end{note}

\paragraph{Markley 2003 的误差表示对比。}
Markley 系统比较了旋转向量、Gibbs 向量、MRP、$2\delta\mathbf{q}_v$ 四种三轴姿态误差\cite{markley2003attitude}：一阶下等价，二阶偏置不同，其中\emph{旋转向量与 MRP 数值性能最佳、适用范围最广}。同文引入的“范数保持四元数复位”正是 \cref{ch:eskf} 复位步（\cref{eq:eskf-reset}）的理论源头。

\begin{table}[H]\centering\small
\caption{MEKF 与 ESKF：特例与推广}
\label{tab:mekf-eskf}
\begin{tabular}{lll}
\toprule
方面 & MEKF & ESKF \\
\midrule
状态 & 姿态 $+$ 陀螺零偏（$6$ 维） & 姿态 $+$ 位置 $+$ 速度 $+$ 双零偏 $+$ 重力（$15$--$18$ 维） \\
位置/速度 & 无 & 加性误差 \\
姿态误差 & \textbf{乘性} $\delta\mathbf{q}$ & \textbf{乘性} $\delta\mathbf{q}$（继承自 MEKF） \\
应用 & 航天器 AOCS、手机 AHRS & VIO、惯性导航、LiDAR-IO \\
约定 & 常用 JPL & Hamilton/JPL 均可 \\
\bottomrule
\end{tabular}
\end{table}

\noindent{}一句话：\textbf{ESKF $=$ MEKF 的 $\SE(3)$/IMU 驱动推广}。两者共享“乘性误差 $+$ 最小维度 $+$ 零点线性化”三大内核；现代视角（Solà 微李理论\cite{sola2018micro}）用 $\boxplus/\boxminus$（\cref{eq:boxminus-rot}）把二者统一。

\subsection{从 MEKF 到 ESKF：姿态误差动力学的逐步推导}
\label{subsec:mekf-to-eskf}

\cref{ch:eskf} 的 \cref{thm:eskf-error-kin} 直接给出了误差动力学结论。这里给\emph{自包含}的逐步推导，让读者看清每条 $\mathbf{F}$ 块从何而来。先看最小问题：只估计姿态与陀螺零偏。

\paragraph{建模。}陀螺测量模型 $\boldsymbol{\omega}_m=\boldsymbol{\omega}+\mathbf{b}_g+\mathbf{n}_g$（真角速度 $+$ 零偏 $+$ 白噪声）。真姿态与标称姿态分别满足
\begin{equation}
  \dot{\mathbf{R}}=\mathbf{R}(\boldsymbol{\omega}_m-\mathbf{b}_g-\mathbf{n}_g)^\wedge,\qquad
  \dot{\hat{\mathbf{R}}}=\hat{\mathbf{R}}(\boldsymbol{\omega}_m-\hat{\mathbf{b}}_g)^\wedge,
  \label{eq:mekf-true-nominal}
\end{equation}
标称态不含瞬时噪声、只用\emph{估计}零偏积分。定义\emph{右乘}误差（与本书右扰动主线一致）
\begin{equation}
  \mathbf{R}=\hat{\mathbf{R}}\,\mathrm{Exp}(\delta\boldsymbol{\theta}),
  \label{eq:mekf-err-def}
\end{equation}
动机：$\hat{\mathbf{R}}$ 表示大姿态，$\mathrm{Exp}(\delta\boldsymbol{\theta})$ 表示从标称到真值的\emph{小转动}，小角误差 $\delta\boldsymbol{\theta}$ 才是 KF 要估的量。

\begin{derivation}[MEKF/ESKF 角误差动力学]
\label{deriv:mekf-theta}
\textbf{Step 1（对复合关系求导）.}\; 对 \cref{eq:mekf-err-def} 求导，并用小角近似 $\mathrm{Exp}(\delta\boldsymbol{\theta})\approx\mathbf{I}+\delta\boldsymbol{\theta}^\wedge$、$\tfrac{d}{dt}\mathrm{Exp}(\delta\boldsymbol{\theta})\approx\dot{\delta\boldsymbol{\theta}}^\wedge$：
\[
  \dot{\mathbf{R}}=\dot{\hat{\mathbf{R}}}\,\mathrm{Exp}(\delta\boldsymbol{\theta})+\hat{\mathbf{R}}\tfrac{d}{dt}\mathrm{Exp}(\delta\boldsymbol{\theta})
  \approx\hat{\mathbf{R}}(\boldsymbol{\omega}_m-\hat{\mathbf{b}}_g)^\wedge(\mathbf{I}+\delta\boldsymbol{\theta}^\wedge)+\hat{\mathbf{R}}\dot{\delta\boldsymbol{\theta}}^\wedge.
\]
记 $\hat{\boldsymbol{\omega}}=\boldsymbol{\omega}_m-\hat{\mathbf{b}}_g$。左乘 $\hat{\mathbf{R}}^\top$，左边一阶项为 $\hat{\boldsymbol{\omega}}^\wedge+\hat{\boldsymbol{\omega}}^\wedge\delta\boldsymbol{\theta}^\wedge+\dot{\delta\boldsymbol{\theta}}^\wedge$。

\textbf{Step 2（真实动力学写到标称坐标）.}\; 把真实 \cref{eq:mekf-true-nominal} 也代入 \cref{eq:mekf-err-def}，记 $\delta\mathbf{b}_g=\mathbf{b}_g-\hat{\mathbf{b}}_g$，则 $\boldsymbol{\omega}_m-\mathbf{b}_g-\mathbf{n}_g=\hat{\boldsymbol{\omega}}-\delta\mathbf{b}_g-\mathbf{n}_g$，忽略二阶小量 $\delta\boldsymbol{\theta}\,\delta\mathbf{b}_g$、$\delta\boldsymbol{\theta}\,\mathbf{n}_g$，左乘 $\hat{\mathbf{R}}^\top$ 后右边一阶项为 $\hat{\boldsymbol{\omega}}^\wedge+\delta\boldsymbol{\theta}^\wedge\hat{\boldsymbol{\omega}}^\wedge-(\delta\mathbf{b}_g+\mathbf{n}_g)^\wedge$。

\textbf{Step 3（比较一阶项）.}\; 消去公共的 $\hat{\boldsymbol{\omega}}^\wedge$：
\[
  \dot{\delta\boldsymbol{\theta}}^\wedge=\delta\boldsymbol{\theta}^\wedge\hat{\boldsymbol{\omega}}^\wedge-\hat{\boldsymbol{\omega}}^\wedge\delta\boldsymbol{\theta}^\wedge-(\delta\mathbf{b}_g+\mathbf{n}_g)^\wedge.
\]
用李括号恒等式 $\mathbf{a}^\wedge\mathbf{b}^\wedge-\mathbf{b}^\wedge\mathbf{a}^\wedge=(\mathbf{a}\times\mathbf{b})^\wedge$（\cref{ch:lie}），得 $\delta\boldsymbol{\theta}^\wedge\hat{\boldsymbol{\omega}}^\wedge-\hat{\boldsymbol{\omega}}^\wedge\delta\boldsymbol{\theta}^\wedge=(\delta\boldsymbol{\theta}\times\hat{\boldsymbol{\omega}})^\wedge=(-\hat{\boldsymbol{\omega}}^\wedge\delta\boldsymbol{\theta})^\wedge$，于是
\begin{equation}
  \boxed{\;\dot{\delta\boldsymbol{\theta}}=-\hat{\boldsymbol{\omega}}^\wedge\delta\boldsymbol{\theta}-\delta\mathbf{b}_g-\mathbf{n}_g.\;}
  \label{eq:mekf-theta-dyn}
\end{equation}
\end{derivation}
\cref{eq:mekf-theta-dyn} 即 MEKF 姿态误差方程，也是 \cref{ch:eskf} 的 \cref{thm:eskf-error-kin} 中角误差块 $\dot{\delta\boldsymbol{\theta}}=-[\boldsymbol{\omega}_m-\boldsymbol{\omega}_b]_\times\delta\boldsymbol{\theta}-\delta\boldsymbol{\omega}_b-\boldsymbol{\omega}_n$ 的来源。

\paragraph{推广到 ESKF 的速度误差。}ESKF 把状态扩成 $(\mathbf{p},\mathbf{v},\mathbf{R},\mathbf{b}_a,\mathbf{b}_g,\mathbf{g})$，复合关系姿态乘性、其余加性（\cref{ch:eskf} 的 \cref{def:eskf-states}）。从 IMU 比力模型 $\dot{\mathbf{v}}=\mathbf{g}+\mathbf{R}(\mathbf{a}_m-\mathbf{b}_a-\mathbf{n}_a)$ 出发，记 $\hat{\mathbf{a}}=\mathbf{a}_m-\hat{\mathbf{b}}_a$，一阶展开
\[
  \mathbf{R}(\mathbf{a}_m-\mathbf{b}_a-\mathbf{n}_a)\approx\hat{\mathbf{R}}(\mathbf{I}+\delta\boldsymbol{\theta}^\wedge)(\hat{\mathbf{a}}-\delta\mathbf{b}_a-\mathbf{n}_a)
  \approx\hat{\mathbf{R}}\hat{\mathbf{a}}+\hat{\mathbf{R}}\delta\boldsymbol{\theta}^\wedge\hat{\mathbf{a}}-\hat{\mathbf{R}}\delta\mathbf{b}_a-\hat{\mathbf{R}}\mathbf{n}_a,
\]
由 $\delta\boldsymbol{\theta}^\wedge\hat{\mathbf{a}}=-\hat{\mathbf{a}}^\wedge\delta\boldsymbol{\theta}$ 得速度误差与位置误差
\begin{equation}
  \dot{\delta\mathbf{v}}=-\hat{\mathbf{R}}\hat{\mathbf{a}}^\wedge\delta\boldsymbol{\theta}-\hat{\mathbf{R}}\delta\mathbf{b}_a+\delta\mathbf{g}-\hat{\mathbf{R}}\mathbf{n}_a,\qquad
  \dot{\delta\mathbf{p}}=\delta\mathbf{v},
  \label{eq:eskf-v-dyn}
\end{equation}
与 \cref{ch:eskf} 的 \cref{thm:eskf-error-kin} 一致。可见 \cref{ch:eskf} 的 $\mathbf{F}$ 矩阵不是凭空背出来的，而是由三个复合关系逐块推导（角误差见 \cref{deriv:mekf-theta}、速度/位置见 \cref{eq:eskf-v-dyn}）。

\begin{pitfall}[把 $\delta\boldsymbol{\theta}$ 当欧拉角]
$\delta\boldsymbol{\theta}$ 是\emph{切空间坐标}（单位元附近的小转动向量），不是全局姿态参数。大姿态机动后若不每步把误差注入标称态并\emph{复位}（\cref{eq:eskf-reset}），$\delta\boldsymbol{\theta}$ 会越积越大、线性化失效。这与 \cref{ch:eskf} 强调“误差恒在原点附近”是同一回事。
\end{pitfall}

\begin{practice}
\begin{enumerate}
  \item 独立复现 \cref{deriv:mekf-theta}，指出哪一步丢了二阶项、为什么可以丢。
  \item 把 \cref{eq:mekf-err-def} 改成\emph{左乘}误差 $\mathbf{R}=\mathrm{Exp}(\delta\boldsymbol{\theta})\hat{\mathbf{R}}$ 重做 Step 1--3，对比 \cref{eq:mekf-theta-dyn} 哪些符号变（这正是 \cref{ch:eskf} 的 \cref{tab:eskf-pert} 左右扰动对照的来由）。
  \item 由 \cref{eq:eskf-v-dyn} 说明：为什么速度误差里出现 $-\hat{\mathbf{R}}\hat{\mathbf{a}}^\wedge\delta\boldsymbol{\theta}$ 这一“姿态—速度耦合”项？它在 \cref{sec:log-linear} 的 InEKF 里会变成什么？
\end{enumerate}
\end{practice}

% ════════════════════════════════════════════════════════════════
\section{群仿射系统：不变误差为何自治}
\label{sec:group-affine}

\paragraph{动机：ESKF 改善了一致性，但没根治。}
\cref{ch:eskf} 已点明（\cref{paper:inekf}）：ESKF 把误差压在原点附近、改善了一致性，却\emph{没有}消除“线性化雅可比依赖估计状态”这一病根——速度误差里的 $-\hat{\mathbf{R}}\hat{\mathbf{a}}^\wedge\delta\boldsymbol{\theta}$（\cref{eq:eskf-v-dyn}）就含 $\hat{\mathbf{R}}$。Barrau--Bonnabel\cite{barrau2017invariant} 发现了一类\emph{特殊}的非线性系统——群仿射系统——在其上误差动力学\emph{精确}地与状态轨迹解耦。\cref{ch:lie} 的 \cref{sec:lie-symmetry} 已给对称/不变/等变的概念地基；本节给完整代数。

\subsection{左/右不变误差}
\label{subsec:invariant-error}

设状态 $\mathbf{X}_t\in G$（矩阵李群），其动力学
\begin{equation}
  \dot{\mathbf{X}}_t=f_{\mathbf{u}_t}(\mathbf{X}_t),
  \label{eq:ga-dynamics}
\end{equation}
$\mathbf{u}_t$ 为输入。对系统的任意两条轨迹 $\mathbf{X}_t,\hat{\mathbf{X}}_t$（抽象阶段不区分真值与估计；进入滤波阶段则 $\mathbf{X}_t$ 为真值、$\hat{\mathbf{X}}_t$ 为估计），定义两种\emph{不变误差}（与 \cref{def:invariant-error} 一致，这里用群元素而非其对数）：
\begin{equation}
  \boldsymbol{\eta}^L_t:=\mathbf{X}_t^{-1}\hat{\mathbf{X}}_t\quad\text{(左不变)},\qquad
  \boldsymbol{\eta}^R_t:=\hat{\mathbf{X}}_t\mathbf{X}_t^{-1}\quad\text{(右不变)}.
  \label{eq:lr-error}
\end{equation}
“左不变”得名于：$\boldsymbol{\eta}^L$ 在对两条轨迹同时\emph{左乘} $\boldsymbol{\Gamma}\in G$（即 $(\mathbf{X},\hat{\mathbf{X}})\mapsto(\boldsymbol{\Gamma}\mathbf{X},\boldsymbol{\Gamma}\hat{\mathbf{X}})$）下不变；右不变同理。

\begin{note}
\textbf{严重的约定坑（务必读）.}\;
\cref{eq:lr-error} 采用 Barrau--Bonnabel TAC 2017 原始约定（\emph{估计在右}：$\boldsymbol{\eta}^L=\mathbf{X}^{-1}\hat{\mathbf{X}}$）。但 Hartley 等\cite{hartley2020contact}、Bonnabel 旧文、部分 MEKF 文献把右不变误差写成 $\hat{\mathbf{X}}\mathbf{X}^{-1}$、左不变写成 $\mathbf{X}\hat{\mathbf{X}}^{-1}$——交换 $\mathbf{X}\leftrightarrow\hat{\mathbf{X}}$ 即互为逆、恰好对称换位。manif 库的 \texttt{rminus}/\allowbreak\texttt{lminus} 又是另一套命名。\cref{ch:eskf} 的 \cref{paper:inekf} 取右不变 $\boldsymbol{\eta}=\hat{\mathbf{X}}\mathbf{X}^{-1}$（与本节 $\boldsymbol{\eta}^R$ 一致）。\textbf{交叉阅读论文时，必须先以一处定义为基准重新换算，否则所有符号全错。}
\end{note}

\subsection{群仿射条件与三等价定理}
\label{subsec:ga-theorem}

\begin{definition}[群仿射动力学]\label{def:group-affine}
\cref{eq:ga-dynamics} 的动力学称为\emph{群仿射}（group-affine），当且仅当对所有 $t$ 与所有 $\mathbf{a},\mathbf{b}\in G$，
\begin{equation}
  \boxed{\;f_{\mathbf{u}_t}(\mathbf{a}\mathbf{b})=f_{\mathbf{u}_t}(\mathbf{a})\,\mathbf{b}+\mathbf{a}\,f_{\mathbf{u}_t}(\mathbf{b})-\mathbf{a}\,f_{\mathbf{u}_t}(\mathbf{e})\,\mathbf{b},\;}
  \label{eq:group-affine}
\end{equation}
其中 $\mathbf{e}$ 为单位元。
\end{definition}

\begin{insight}[为什么叫“仿射”]
在 $(G,+)=\mathbb{R}^n$（加法群、$\mathbf{e}=\mathbf{0}$）上，\cref{eq:group-affine} 退化为 $f(\mathbf{a}+\mathbf{b})=f(\mathbf{a})+f(\mathbf{b})-f(\mathbf{0})$。令 $g(\mathbf{x}):=f(\mathbf{x})-f(\mathbf{0})$，则 $g$ 满足可加性 $\Rightarrow$ $g$ 线性 $\Rightarrow$ $f(\mathbf{x})=\mathbf{A}\mathbf{x}+\mathbf{b}$ 即中学意义的\emph{仿射函数}。把它搬到李群上就是 group-affine。关键：$f(\mathbf{e})$ \emph{不必为零}——这是“仿射”而非“线性”的本质来源，重力 $\mathbf{g}$ 即 $f(\mathbf{e})$ 的非零项（见 \cref{subsec:se23-verify}）。
\end{insight}

\begin{theorem}[群仿射三等价（Barrau--Bonnabel TAC 2017, Thm 1）]\label{thm:ga-equiv}
对动力学 \cref{eq:ga-dynamics}，下列三条件\emph{等价}：
\begin{enumerate}
  \item[(i)] 左不变误差具状态轨迹无关传播：存在 $g^L_{\mathbf{u}_t}$ 使 $\dot{\boldsymbol{\eta}}^L_t=g^L_{\mathbf{u}_t}(\boldsymbol{\eta}^L_t)$；
  \item[(ii)] 右不变误差具同样性质：$\dot{\boldsymbol{\eta}}^R_t=g^R_{\mathbf{u}_t}(\boldsymbol{\eta}^R_t)$；
  \item[(iii)] 群仿射条件 \cref{eq:group-affine} 成立。
\end{enumerate}
此时（取右不变误差，估计在右 $\boldsymbol{\eta}=\hat{\mathbf{X}}\mathbf{X}^{-1}$）误差动力学为
\begin{equation}
  g^L_{\mathbf{u}_t}(\boldsymbol{\eta})=f_{\mathbf{u}_t}(\boldsymbol{\eta})-f_{\mathbf{u}_t}(\mathbf{e})\,\boldsymbol{\eta},\qquad
  g^R_{\mathbf{u}_t}(\boldsymbol{\eta})=f_{\mathbf{u}_t}(\boldsymbol{\eta})-\boldsymbol{\eta}\,f_{\mathbf{u}_t}(\mathbf{e}).
  \label{eq:ga-error-dyn}
\end{equation}
\end{theorem}

\begin{derivation}[\cref{thm:ga-equiv} 的完整证明]
\label{deriv:ga-equiv}
\textbf{预备引理.}\; 对任意可微曲线 $\mathbf{X}_t\in G$，由 $\mathbf{X}_t\mathbf{X}_t^{-1}=\mathbf{I}$ 求导得
\begin{equation}
  \tfrac{d}{dt}(\mathbf{X}_t^{-1})=-\mathbf{X}_t^{-1}\dot{\mathbf{X}}_t\mathbf{X}_t^{-1}.
  \label{eq:inv-deriv}
\end{equation}

\textbf{(i) $\Rightarrow$ (iii).}\; 设 $\mathbf{X}_t,\hat{\mathbf{X}}_t$ 为两条轨迹，$\boldsymbol{\eta}_t=\mathbf{X}_t^{-1}\hat{\mathbf{X}}_t$（省略上标）。由 \cref{eq:inv-deriv}、代入动力学、并用 $\mathbf{X}_t^{-1}\hat{\mathbf{X}}_t=\boldsymbol{\eta}_t$、$\hat{\mathbf{X}}_t=\mathbf{X}_t\boldsymbol{\eta}_t$：
\begin{equation}
  \dot{\boldsymbol{\eta}}_t=-\mathbf{X}_t^{-1}f_{\mathbf{u}_t}(\mathbf{X}_t)\boldsymbol{\eta}_t+\mathbf{X}_t^{-1}f_{\mathbf{u}_t}(\mathbf{X}_t\boldsymbol{\eta}_t).
  \label{eq:ga-eta-exact}
\end{equation}
此为精确等式（尚未用 (i)）。假设 (i) 成立，\cref{eq:ga-eta-exact} 对所有 $\mathbf{X}_t,\boldsymbol{\eta}_t\in G$ 成立。取特殊值 $\mathbf{X}_t=\mathbf{e}$：$g_{\mathbf{u}_t}(\boldsymbol{\eta}_t)=-f_{\mathbf{u}_t}(\mathbf{e})\boldsymbol{\eta}_t+f_{\mathbf{u}_t}(\boldsymbol{\eta}_t)$，即 \cref{eq:ga-error-dyn} 左式。代回 \cref{eq:ga-eta-exact}，把一般 $\mathbf{X}_t$ 写作 $\mathbf{a}$、$\boldsymbol{\eta}_t$ 写作 $\mathbf{b}$，左乘 $\mathbf{a}$ 整理即得 \cref{eq:group-affine}。

\textbf{(iii) $\Rightarrow$ (i).}\; 由 \cref{eq:group-affine} 取 $\mathbf{a}=\mathbf{X}_t,\mathbf{b}=\boldsymbol{\eta}_t$，代入 \cref{eq:ga-eta-exact}：$-\mathbf{X}_t^{-1}f_{\mathbf{u}_t}(\mathbf{X}_t)\boldsymbol{\eta}_t$ 与展开后的 $+\mathbf{X}_t^{-1}f_{\mathbf{u}_t}(\mathbf{X}_t)\boldsymbol{\eta}_t$ 互消，剩 $f_{\mathbf{u}_t}(\boldsymbol{\eta}_t)-f_{\mathbf{u}_t}(\mathbf{e})\boldsymbol{\eta}_t=g^L_{\mathbf{u}_t}(\boldsymbol{\eta}_t)$，$\mathbf{X}_t$ \emph{完全消除}，(i) 得证。

\textbf{(ii) $\Leftrightarrow$ (iii).}\; 对 $\tilde{\boldsymbol{\eta}}_t=\hat{\mathbf{X}}_t\mathbf{X}_t^{-1}$ 同法求导、取 $\mathbf{X}_t=\mathbf{e}$ 得 \cref{eq:ga-error-dyn} 右式；用 \cref{eq:group-affine} 取 $\mathbf{a}=\tilde{\boldsymbol{\eta}}_t,\mathbf{b}=\mathbf{X}_t$ 平行展开即得，$\mathbf{X}_t$ 消除。$\square$
\end{derivation}

\begin{insight}[$f(\mathbf{e})$ 的“常数漂移”角色]
若 $f_{\mathbf{u}}(\mathbf{e})=\mathbf{0}$，群仿射退化为左/右不变的纯乘法形式（$f(\mathbf{X})=\mathbf{X}\boldsymbol{\omega}$ 纯左不变、$f(\mathbf{X})=\mathbf{v}\mathbf{X}$ 纯右不变都满足 \cref{eq:group-affine}）。一般情况下 $f_{\mathbf{u}}(\mathbf{e})$ 扮演“常数项”，类比 $\mathbb{R}^n$ 上仿射函数的 $f(\mathbf{0})$。\emph{group-affine 严格扩张了左右不变的并集}——下节将看到 $\mathrm{SE}_2(3)$ 上 IMU 动力学既非左不变也非右不变，但仍群仿射，秘密就在重力使 $f(\mathbf{e})\neq\mathbf{0}$。
\end{insight}

\begin{pitfall}[概念误区：“group-affine 就是左不变加右不变”]
初学者常把群仿射等同于“左不变 $+$ 右不变”。错。区分三者的判据是看 $f(\mathbf{e})$ 是否为零：$f(\mathbf{e})=\mathbf{0}$ 时退化为左/右不变；$f(\mathbf{e})\neq\mathbf{0}$ 时“仿射”项登场，这才是 group-affine 的\emph{真正}新意（也是 Barrau 博士论文的核心发现：以前以为只有左/右不变系统能做不变滤波，IMU 导航两者都不是，但它\emph{是}群仿射的）。
\end{pitfall}

\subsection{\texorpdfstring{$\mathrm{SE}_2(3)$}{SE2(3)} 上 IMU 动力学的逐块验证}
\label{subsec:se23-verify}

\cref{ch:lie} 的 \cref{sec:se23} 已给出扩展位姿群 $\mathrm{SE}_2(3)$ 的权威定义（\cref{def:se23}）、指数映射（\cref{thm:se23-exp}）与伴随（\cref{eq:se23-adjoint}），并在 \cref{sec:se23} 末已指出“IMU 运动学在 $\mathrm{SE}_2(3)$ 上群仿射”。本节给\emph{逐块}验证——很多工程错误恰恰发生在“知道它成立、却不知每块为何成立”的断层上。

\paragraph{设定。}IMU 连续模型（陀螺 $\boldsymbol{\omega}$、加速度计 $\mathbf{a}$、重力 $\mathbf{g}$）
\begin{equation}
  \dot{\mathbf{R}}=\mathbf{R}\boldsymbol{\omega}^\wedge,\qquad \dot{\mathbf{v}}=\mathbf{g}+\mathbf{R}\mathbf{a},\qquad \dot{\mathbf{r}}=\mathbf{v}.
  \label{eq:imu-kin}
\end{equation}
按 \cref{def:se23} 写成 $5\times5$ 矩阵 $\mathbf{X}=\big[\begin{smallmatrix}\mathbf{R}&\mathbf{r}&\mathbf{v}\\\mathbf{0}^\top&1&0\\\mathbf{0}^\top&0&1\end{smallmatrix}\big]$（位置列 $\mathbf{r}$ 在前、速度列 $\mathbf{v}$ 在后），动力学的矩阵形式为
\begin{equation}
  f_{\boldsymbol{\omega},\mathbf{a}}(\mathbf{X})=\begin{bmatrix}\mathbf{R}\boldsymbol{\omega}^\wedge&\mathbf{v}&\mathbf{g}+\mathbf{R}\mathbf{a}\\\mathbf{0}^\top&0&0\\\mathbf{0}^\top&0&0\end{bmatrix},\qquad
  f_{\boldsymbol{\omega},\mathbf{a}}(\mathbf{I})=\begin{bmatrix}\boldsymbol{\omega}^\wedge&\mathbf{0}&\mathbf{g}+\mathbf{a}\\\mathbf{0}^\top&0&0\\\mathbf{0}^\top&0&0\end{bmatrix}.
  \label{eq:imu-matrix-f}
\end{equation}

\begin{note}
\textbf{排序对照.}\;
源文献\cite{barrau2017invariant,hartley2020contact} 把\emph{速度列放位置列之前}、且李代数\emph{旋转在最前}，其 $f(\mathbf{X})$ 第二列是 $\mathbf{g}+\mathbf{R}\mathbf{a}$、第三列是 $\mathbf{v}$。\cref{eq:imu-matrix-f} 已按本书 \cref{def:se23} 的“位置列在前、速度列在后”转写（第二列 $\mathbf{v}$、第三列 $\mathbf{g}+\mathbf{R}\mathbf{a}$）。两者只是列的置换，群仿射性与下面验证不受影响。
\end{note}

\begin{derivation}[$\mathrm{SE}_2(3)$ IMU 动力学逐块满足 \cref{eq:group-affine}]
\label{deriv:se23-ga}
最直接的看法：\cref{eq:imu-matrix-f} 可写成 $f(\mathbf{X})=\mathbf{A}\mathbf{X}+\mathbf{X}\mathbf{B}$，其中 $\mathbf{A}$ 携带重力（$f(\mathbf{e})$ 的非零部分）、$\mathbf{B}$ 携带机体系输入。对任意 $\mathbf{X}_1,\mathbf{X}_2$，
\[
  f(\mathbf{X}_1)\mathbf{X}_2+\mathbf{X}_1 f(\mathbf{X}_2)-\mathbf{X}_1 f(\mathbf{I})\mathbf{X}_2
  =(\mathbf{A}\mathbf{X}_1+\mathbf{X}_1\mathbf{B})\mathbf{X}_2+\mathbf{X}_1(\mathbf{A}\mathbf{X}_2+\mathbf{X}_2\mathbf{B})-\mathbf{X}_1(\mathbf{A}+\mathbf{B})\mathbf{X}_2
  =\mathbf{A}\mathbf{X}_1\mathbf{X}_2+\mathbf{X}_1\mathbf{X}_2\mathbf{B}=f(\mathbf{X}_1\mathbf{X}_2),
\]
其中 $\mathbf{X}_1\mathbf{B}\mathbf{X}_2$ 与 $\mathbf{X}_1\mathbf{A}\mathbf{X}_2$ 各被加一次又减一次而抵消。这就是 \cref{eq:group-affine}（$f(\mathbf{e})=\mathbf{A}+\mathbf{B}$）。\emph{逐块}核对（群乘法 $\mathbf{X}_1\mathbf{X}_2=\big[\begin{smallmatrix}\mathbf{R}_1\mathbf{R}_2&\mathbf{R}_1\mathbf{r}_2+\mathbf{r}_1&\mathbf{R}_1\mathbf{v}_2+\mathbf{v}_1\\&1&\\&&1\end{smallmatrix}\big]$，\cref{prop:se23-group}）：

\textbf{旋转块.}\; 左边 $\mathbf{R}_1\mathbf{R}_2\boldsymbol{\omega}^\wedge$；右边 $\mathbf{R}_1\boldsymbol{\omega}^\wedge\mathbf{R}_2+\mathbf{R}_1\mathbf{R}_2\boldsymbol{\omega}^\wedge-\mathbf{R}_1\boldsymbol{\omega}^\wedge\mathbf{R}_2=\mathbf{R}_1\mathbf{R}_2\boldsymbol{\omega}^\wedge$。$\checkmark$（抵消项正是 $-\mathbf{X}_1 f(\mathbf{I})\mathbf{X}_2$ 的作用）。

\textbf{速度列.}\; 左边 $\mathbf{g}+\mathbf{R}_1\mathbf{R}_2\mathbf{a}$；右边三项的速度列依次为 $\mathbf{R}_1\boldsymbol{\omega}^\wedge\mathbf{v}_2+\mathbf{g}+\mathbf{R}_1\mathbf{a}$、$\mathbf{R}_1(\mathbf{g}+\mathbf{R}_2\mathbf{a})$、$\mathbf{R}_1\boldsymbol{\omega}^\wedge\mathbf{v}_2+\mathbf{R}_1(\mathbf{g}+\mathbf{a})$，合并得 $\mathbf{g}+\mathbf{R}_1\mathbf{R}_2\mathbf{a}$。$\checkmark$（这里最能看出“仿射”：世界系重力 $\mathbf{g}$ 不随 $\mathbf{R}_1$ 旋转，减项正好消掉多余的 $\mathbf{R}_1\mathbf{g}$）。

\textbf{位置列.}\; 左边 $\mathbf{R}_1\mathbf{v}_2+\mathbf{v}_1$；右边三项位置列 $\mathbf{R}_1\boldsymbol{\omega}^\wedge\mathbf{r}_2+\mathbf{v}_1$、$\mathbf{R}_1\mathbf{v}_2$、$\mathbf{R}_1\boldsymbol{\omega}^\wedge\mathbf{r}_2$，合并得 $\mathbf{R}_1\mathbf{v}_2+\mathbf{v}_1$。$\checkmark$ 三块全部验证完毕。
\end{derivation}

\begin{insight}[$\mathrm{SE}_2(3)$ 为何不可或缺]
逐块验证也说明普通 $\SE(3)$ \emph{不够}：$\SE(3)$ 只含 $\mathbf{R},\mathbf{r}$，速度 $\mathbf{v}$ 在群外，$\dot{\mathbf{r}}=\mathbf{v}$ 不能作为群乘法的一列自然出现。$\mathrm{SE}_2(3)$ 把速度并入矩阵（\cref{def:se23}），正是为惯导量身定制——\emph{把 IMU 动力学变成 $\mathbf{A}\mathbf{X}+\mathbf{X}\mathbf{B}$ 的代数形态}，群仿射几乎自动成立。这与 \cref{ch:lie} 的 \cref{sec:se23} “为何要费这个劲”遥相呼应。
\end{insight}

\begin{practice}
\begin{enumerate}
  \item 验证纯左不变 $f(\mathbf{X})=\mathbf{X}\boldsymbol{\omega}$ 与纯右不变 $f(\mathbf{X})=\mathbf{v}\mathbf{X}$ 都满足 \cref{eq:group-affine}（提示：直接代入，$f(\mathbf{e})$ 分别为 $\boldsymbol{\omega},\mathbf{v}$）。
  \item 在 \cref{deriv:se23-ga} 中去掉重力 $\mathbf{g}$，验证此时 $f(\mathbf{I})$ 如何变化、系统为何\emph{仍}群仿射。
  \item 用 \cref{thm:ga-equiv} 的 \cref{eq:ga-error-dyn}（右不变）写出 $\mathrm{SE}_2(3)$ IMU 的误差动力学 $\dot{\boldsymbol{\eta}}^R=f(\boldsymbol{\eta}^R)-\boldsymbol{\eta}^R f(\mathbf{e})$，确认其右端\emph{不含}真实状态 $\mathbf{X}$。
\end{enumerate}
\end{practice}

% ════════════════════════════════════════════════════════════════
\section{log-linear 定理：误差对数的精确线性}
\label{sec:log-linear}

\paragraph{动机：自治还不够，我们要“线性”。}
\cref{thm:ga-equiv} 给出误差动力学 $\dot{\boldsymbol{\eta}}=g_{\mathbf{u}_t}(\boldsymbol{\eta})$ 自治（不依赖状态）。但 $g_{\mathbf{u}_t}$ 在群上仍\emph{非线性}。Barrau--Bonnabel 的第二个奇迹是：把误差取对数坐标 $\boldsymbol{\eta}_t=\mathrm{Exp}(\boldsymbol{\xi}_t)$ 后，$\boldsymbol{\xi}_t$ 满足\emph{精确}的线性时变（LTV）方程——而非通常的一阶近似。这正是 \cref{ch:eskf} 的 \cref{paper:inekf} 所称 “log-linear property”，本节给其机理。

\begin{theorem}[log-linear 性质（Barrau--Bonnabel TAC 2017, Thm 2）]\label{thm:log-linear}
设 \cref{eq:ga-dynamics} 群仿射，$\boldsymbol{\eta}^i_t$（$i\in\{L,R\}$）为不变误差，取 $\boldsymbol{\xi}^i_0$ 使 $\boldsymbol{\eta}^i_0=\mathrm{Exp}(\boldsymbol{\xi}^i_0)$。设矩阵 $\mathbf{A}^i_{\mathbf{u}_t}$ 由 $g^i_{\mathbf{u}_t}$ 在单位元的一阶展开 $g^i_{\mathbf{u}_t}(\mathrm{Exp}(\boldsymbol{\xi}))=\big(\mathbf{A}^i_{\mathbf{u}_t}\boldsymbol{\xi}\big)^\wedge+\mathcal{O}(\|\boldsymbol{\xi}\|^2)$ 定义。若 $\boldsymbol{\xi}^i_t$ 满足线性 ODE
\begin{equation}
  \dot{\boldsymbol{\xi}}^i_t=\mathbf{A}^i_{\mathbf{u}_t}\boldsymbol{\xi}^i_t,
  \label{eq:log-linear}
\end{equation}
则对\emph{任意} $t\ge0$（且初始误差只需可由所选 $\log/\exp$ 分支表示，\emph{不要求小}）都有 $\boldsymbol{\eta}^i_t=\mathrm{Exp}(\boldsymbol{\xi}^i_t)$。即 \cref{eq:log-linear} 是\textbf{精确}描述（无噪声）非线性误差传播的方程，$\mathbf{A}^i_{\mathbf{u}_t}$ \emph{只依赖输入} $\mathbf{u}_t$、与 $\mathbf{X}_t,\hat{\mathbf{X}}_t$ 均无关。
\end{theorem}

\begin{derivation}[log-linear 的 BCH 机理]
\label{deriv:log-linear}
核心是证明误差流 $\boldsymbol{\Phi}_t$（即 $\boldsymbol{\eta}_t=\boldsymbol{\Phi}_t(\boldsymbol{\eta}_0)$）是群\emph{同态}，再用 BCH 平凡化。

\textbf{Part A（流的存在）.}\; 由 \cref{thm:ga-equiv}，群仿射 $\Leftrightarrow$ $\dot{\boldsymbol{\eta}}_t=g^L_{\mathbf{u}_t}(\boldsymbol{\eta}_t)$ 右端只依赖 $\boldsymbol{\eta}_t$。$g^L_{\mathbf{u}_t}$ 局部 Lipschitz（由 $f_{\mathbf{u}}$ 光滑保证），Picard--Lindelöf 给出唯一流 $\boldsymbol{\Phi}_t:G\to G$，仅依赖输入 $\{\mathbf{u}_s\}$。

\textbf{Part B（流是群同态，关键引理）.}\; $g^L_{\mathbf{u}}$ 满足 Leibniz 性质（由 \cref{eq:group-affine} 与 \cref{eq:ga-error-dyn} 推出）：
\begin{equation}
  g^L_{\mathbf{u}}(\mathbf{a}\mathbf{b})=f_{\mathbf{u}}(\mathbf{a}\mathbf{b})-f_{\mathbf{u}}(\mathbf{e})\mathbf{a}\mathbf{b}
  \overset{\cref{eq:group-affine}}{=}g^L_{\mathbf{u}}(\mathbf{a})\mathbf{b}+\mathbf{a}\,g^L_{\mathbf{u}}(\mathbf{b}).
  \label{eq:leibniz}
\end{equation}
对 $\boldsymbol{\Phi}_t(\boldsymbol{\eta}_0)\boldsymbol{\Phi}_t(\boldsymbol{\eta}_0')$ 求导，用 \cref{eq:leibniz} 知其满足与 $\boldsymbol{\Phi}_t(\boldsymbol{\eta}_0\boldsymbol{\eta}_0')$ 相同的 ODE、相同初值，由解唯一性
\begin{equation}
  \boldsymbol{\Phi}_t(\boldsymbol{\eta}_0\boldsymbol{\eta}_0')=\boldsymbol{\Phi}_t(\boldsymbol{\eta}_0)\,\boldsymbol{\Phi}_t(\boldsymbol{\eta}_0').
  \label{eq:flow-homo}
\end{equation}

\textbf{Part C（BCH 平凡化）.}\; 由 \cref{eq:flow-homo} 归纳，$\boldsymbol{\Phi}_t(\boldsymbol{\eta}_0^p)=\boldsymbol{\Phi}_t(\boldsymbol{\eta}_0)^p$。对任意 $n$：$\boldsymbol{\Phi}_t(\mathrm{Exp}(\boldsymbol{\xi}_0))=\boldsymbol{\Phi}_t\big(\mathrm{Exp}(\boldsymbol{\xi}_0/n)\big)^n$。设 $\mathbf{F}_t:=D\boldsymbol{\Phi}_t|_{\mathbf{e}}$，泰勒展开 $\boldsymbol{\Phi}_t(\mathrm{Exp}(\boldsymbol{\xi}_0/n))=\mathrm{Exp}\big(\tfrac1n\mathbf{F}_t\boldsymbol{\xi}_0+\mathbf{r}_t(\boldsymbol{\xi}_0/n)\big)$，$\mathbf{r}_t=\mathcal{O}(\|\cdot\|^2)$。\emph{关键}：$n$ 个\emph{相同}因子 $\mathbf{Y}_n=\tfrac1n\mathbf{F}_t\boldsymbol{\xi}_0+\mathbf{r}_t(\boldsymbol{\xi}_0/n)$ 的 $\mathrm{Exp}$ 乘积，BCH 公式 $\ln(e^{\mathbf{X}}e^{\mathbf{Y}})=\mathbf{X}+\mathbf{Y}+\tfrac12[\mathbf{X},\mathbf{Y}]+\cdots$ 的所有交叉项都含 $[\mathbf{Y}_n,\mathbf{Y}_n]=\mathbf{0}$，故 $[\mathrm{Exp}(\mathbf{Y}_n)]^n=\mathrm{Exp}(n\mathbf{Y}_n)=\mathrm{Exp}\big(\mathbf{F}_t\boldsymbol{\xi}_0+n\,\mathbf{r}_t(\boldsymbol{\xi}_0/n)\big)$。又 $n\,\mathbf{r}_t(\boldsymbol{\xi}_0/n)=n\cdot\mathcal{O}(1/n^2)\to\mathbf{0}$，故
\begin{equation}
  \boldsymbol{\Phi}_t(\mathrm{Exp}(\boldsymbol{\xi}_0))=\mathrm{Exp}(\mathbf{F}_t\boldsymbol{\xi}_0).
  \label{eq:flow-exp-linear}
\end{equation}
对时间求导并在 $\boldsymbol{\xi}_0\to\mathbf{0}$ 比较一阶：$\dot{\mathbf{F}}_t=\mathbf{A}_{\mathbf{u}_t}\mathbf{F}_t$，$\mathbf{F}_0=\mathbf{I}$，即 $\boldsymbol{\xi}_t=\mathbf{F}_t\boldsymbol{\xi}_0$ 满足 \cref{eq:log-linear}。$\square$
\end{derivation}

\begin{insight}[“魔术”全在 $[\mathbf{X},\mathbf{X}]=\mathbf{0}$]
log-linear 的精确性\emph{不靠}小误差近似，而靠两件代数事实：(1) 群仿射使误差流 $\boldsymbol{\Phi}_t$ 成为群同态（\cref{eq:flow-homo}）；(2) 群同态由其在单位元的导数 $\mathbf{A}_t$ \emph{完全决定}，而“$n$ 个相同因子的 $\mathrm{Exp}$ 乘积 $=\mathrm{Exp}$ 之 $n$ 倍”（\cref{eq:flow-exp-linear}）因 $[\mathbf{X},\mathbf{X}]=\mathbf{0}$ 让所有 BCH 交叉项消失。无群仿射，$\boldsymbol{\Phi}_t$ 不是同态、论证崩塌。这就是为什么 group-affine 是\emph{珍贵}的代数条件，而非任意限制。
\end{insight}

\paragraph{$\mathrm{SE}_2(3)$ IMU 的 $\mathbf{A}_t$：严格幂零。}
对 \cref{eq:imu-kin} 的右不变误差，按\emph{源文献序} $\boldsymbol{\xi}=[\delta\boldsymbol{\theta};\delta\mathbf{v};\delta\mathbf{p}]$（旋转—速度—位置，见本章记号约定）算 \cref{eq:log-linear} 的系统矩阵\cite{barrau2017invariant}：
\begin{equation}
  \mathbf{A}=\begin{bmatrix}\mathbf{0}&\mathbf{0}&\mathbf{0}\\\mathbf{g}^\wedge&\mathbf{0}&\mathbf{0}\\\mathbf{0}&\mathbf{I}&\mathbf{0}\end{bmatrix},\qquad
  \mathbf{A}^2=\begin{bmatrix}\mathbf{0}&\mathbf{0}&\mathbf{0}\\\mathbf{0}&\mathbf{0}&\mathbf{0}\\\mathbf{g}^\wedge&\mathbf{0}&\mathbf{0}\end{bmatrix},\qquad \mathbf{A}^3=\mathbf{0},
  \label{eq:se23-A}
\end{equation}
\emph{下三角且严格幂零}（$\mathbf{A}^3=\mathbf{0}$）。于是状态转移矩阵\emph{精确}截断为\cite{barrau2017invariant}
\begin{equation}
  \boldsymbol{\Phi}(\Delta t)=\exp(\mathbf{A}\Delta t)=\mathbf{I}+\mathbf{A}\Delta t+\tfrac12\mathbf{A}^2\Delta t^2
  =\begin{bmatrix}\mathbf{I}&\mathbf{0}&\mathbf{0}\\\mathbf{g}^\wedge\Delta t&\mathbf{I}&\mathbf{0}\\\tfrac12\mathbf{g}^\wedge\Delta t^2&\mathbf{I}\Delta t&\mathbf{I}\end{bmatrix}.
  \label{eq:se23-Phi}
\end{equation}
逐块读出：姿态误差\emph{不被}理想陀螺放大（$\delta\boldsymbol{\theta}_{k+1}=\delta\boldsymbol{\theta}_k$）；姿态误差经重力方向错误进入速度（$\delta\mathbf{v}_{k+1}=\delta\mathbf{v}_k+\mathbf{g}^\wedge\delta\boldsymbol{\theta}_k\Delta t$）；速度误差再进入位置。这与经典惯导误差方程一致，但此处由\emph{李群误差 $+$ 群仿射定理}严格推出。

\begin{note}
\textbf{“精确”的边界（务必澄清）.}\;
log-linear 的精确性\emph{仅在无噪声传播段}成立，且受 $\log$ 映射收敛域限制——$\SO(3)$ 旋转误差接近 $\pi$ 时 $\log$ 多值，BCH 收敛要求 $\|\boldsymbol{\xi}\|$ 在单射半径内。工程滤波里的过程噪声作为一阶随机扰动 $\mathbf{B}_t\mathbf{n}_t$ 叠加进协方差传播（含噪时自治性退化为\emph{近似}，仍优于标准 EKF）。\textbf{切勿}把“无噪声精确线性”误读为“InEKF 全局收敛”或“后验最优”——后者需更强假设，TAC 2017 未声称。另：$\mathfrak{se}_2(3)$ 含旋转部分、整体\emph{不是}幂零李代数，\cref{eq:se23-A} 的 $\mathbf{A}^3=\mathbf{0}$ 是该\emph{特定}右不变误差系统矩阵的幂零性（额外的计算友好性），\emph{不是}李代数幂零；log-linear 精确性本身\emph{不需}幂零，仅需群仿射 $+$ \cref{eq:leibniz}。
\end{note}

\paragraph{这为什么是根本性突破。}对一般非线性 $\dot{\mathbf{x}}=f(\mathbf{x},\mathbf{u})$，标准 EKF 误差 $\mathbf{e}=\mathbf{x}-\hat{\mathbf{x}}$ 满足 $\dot{\mathbf{e}}=\mathbf{A}(\hat{\mathbf{x}},\mathbf{u})\mathbf{e}+\mathcal{O}(\|\mathbf{e}\|^2)$——\emph{雅可比 $\mathbf{A}(\hat{\mathbf{x}},\mathbf{u})$ 依赖估计}。估计漂移 $\to$ $\mathbf{A}$ 偏离真值雅可比 $\to$ Riccati 协方差传播失真 $\to$ Kalman 增益不匹配 $\to$ 估计更漂离 $\to$ 正反馈 $\to$ 滤波器发散。这正是 EKF 稳定性证明必须\emph{假设} $\gamma_{\min}\mathbf{I}\preceq\mathbf{P}_t\preceq\gamma_{\max}\mathbf{I}$ 的根源——而此假设实际\emph{无法}先验验证。InEKF 的 $\mathbf{A}_{\mathbf{u}_t}$ 只依赖输入，Riccati $\dot{\mathbf{P}}=\mathbf{A}_{\mathbf{u}_t}\mathbf{P}+\mathbf{P}\mathbf{A}_{\mathbf{u}_t}^\top+\hat{\mathbf{Q}}_t$ 在传播段\emph{精确}（\cref{thm:bias-overconfident} 是 EKF 过自信的偏置版，这里是其几何根治）。

\begin{table}[H]\centering\small
\caption{标准 EKF vs InEKF：传播雅可比的根本差异}
\label{tab:ekf-vs-inekf}
\begin{tabular}{lll}
\toprule
项目 & 标准 EKF & InEKF（群仿射） \\
\midrule
传播雅可比 & $\mathbf{A}(\hat{\mathbf{x}}_t,\mathbf{u}_t)$ 依赖估计 & $\mathbf{A}(\mathbf{u}_t)$ 仅依赖输入 \\
传播段高阶余项 & $\mathcal{O}(\|\mathbf{e}\|^2)$ 系数随 $\hat{\mathbf{x}}$ 变 & \textbf{为零}（log-linear 精确） \\
一致性修补 & 需 FEJ/OC-EKF（经验） & 不需（结构性，\cref{sec:consistency-inekf}） \\
$\mathbf{P}_t$ 边界 & 需\emph{假设} & 由稳定性定理保证（\cref{sec:stability}） \\
\bottomrule
\end{tabular}
\end{table}

\begin{practice}
\begin{enumerate}
  \item 验证 \cref{eq:se23-A} 的 $\mathbf{A}^3=\mathbf{0}$，并由 \cref{eq:se23-Phi} 写出 $\delta\mathbf{p}_{k+1}$ 的逐块传播式，对照经典惯导误差方程。
  \item 反事实：若把 \cref{eq:se23-Phi} 的精确 $\exp(\mathbf{A}\Delta t)$ 换成一阶欧拉 $\mathbf{I}+\mathbf{A}\Delta t$，在 $\Delta t=0.01\,$s 下 $\tfrac12\mathbf{A}^2\Delta t^2$ 项量级多大？累积 $100$ 步对位置误差的影响（提示：$\sim\tfrac12\|\mathbf{g}\|\Delta t^2$ 量级）。
  \item （概念）用一句话解释：为什么 \cref{deriv:log-linear} 的“魔术”在标准 EKF 上\emph{做不到}（提示：标准 EKF 的误差流不是群同态）。
\end{enumerate}
\end{practice}

% ════════════════════════════════════════════════════════════════
\section{InEKF 算法：右不变滤波的完整流程}
\label{sec:inekf-algo}

\paragraph{动机：把定理落成可运行的 predict-update 循环。}
\cref{sec:group-affine,sec:log-linear} 给了理论；本节给\emph{右不变} InEKF（RIEKF）的完整离散算法。工程中最常见的困惑不是单个公式，而是“这些公式按什么顺序组合、手性怎么配套”。

\subsection{观测分类：左不变 vs 右不变}
\label{subsec:obs-classify}

InEKF 的手性由\emph{观测}决定（这是 \cref{def:invariant-error} 的工程落地）。两类不变观测：
\begin{itemize}
  \item \textbf{左不变观测} $\mathbf{Y}_n=\mathbf{X}_{t_n}\mathbf{d}+\mathbf{V}_n$，$\mathbf{d}$ 是已知的\emph{机体系}参考向量（如 GPS：位置槽选择器 $\mathbf{d}$ 在 body 中已知，被 $\mathbf{X}$ 送到 world）。其创新量与 $\boldsymbol{\eta}^L=\mathbf{X}^{-1}\hat{\mathbf{X}}$ 兼容 $\Rightarrow$ 选 LIEKF。
  \item \textbf{右不变观测} $\mathbf{Y}_n=\mathbf{X}_{t_n}^{-1}\mathbf{d}+\mathbf{V}_n$，$\mathbf{d}$ 是已知的\emph{世界系}参考向量（如 landmark、磁力计、重力、接触足端）。其创新量与 $\boldsymbol{\eta}^R=\hat{\mathbf{X}}\mathbf{X}^{-1}$ 兼容 $\Rightarrow$ 选 RIEKF。
\end{itemize}

\begin{table}[H]\centering\small
\caption{常见传感器的不变观测分类（Barrau--Bonnabel / Hartley）}
\label{tab:obs-classify}
\begin{tabular}{llll}
\toprule
传感器 & 观测方程 & 参考向量坐标系 & 类别 \\
\midrule
GPS（world 位置） & $\mathbf{Y}=\mathbf{X}\mathbf{d}$（$\mathbf{d}$ 选位置槽） & body 固定 & 左不变 \\
Landmark（相机看 world 中 $\mathbf{p}_k$） & $\mathbf{Y}=\mathbf{X}^{-1}\mathbf{p}_k^{\mathrm{hom}}$ & world 固定 & 右不变 \\
接触足端运动学 & $\mathbf{R}^\top(\mathbf{d}-\mathbf{p})$ & world 固定 & 右不变 \\
磁力计（body 读 world $\mathbf{m}_\oplus$） & $\mathbf{Y}=\mathbf{R}^\top\mathbf{m}_\oplus+\mathbf{V}$ & world & 右不变 \\
机体系比力（重力 $\mathbf{g}$） & $\mathbf{Y}=\mathbf{R}^\top\mathbf{g}+\mathbf{V}$ & world & 右不变 \\
\bottomrule
\end{tabular}
\end{table}

\begin{insight}[记忆法则]
参考向量 $\mathbf{d}$ 若“长在 world 里、在 body 里读”（magnetometer、landmark、body velocity）$\Rightarrow$ \emph{右不变}；$\mathbf{d}$ 若“长在 body 里、被 $\mathbf{X}$ 送到 world”（GPS 的 origin 在 body 中已知为位置槽）$\Rightarrow$ \emph{左不变}。最常见的误判是把 landmark 观测 $\mathbf{Y}=\mathbf{R}^\top(\mathbf{p}_k-\mathbf{p})$ 当左不变——它其实是\emph{右不变}，因 $\mathbf{p}_k$ 在 world 中已知（这正是 Barrau--Bonnabel §IV/§V 用 RIEKF 的根源）。
\end{insight}

\subsection{右不变观测的雅可比独立于估计}
\label{subsec:H-independent}

对右不变观测一阶展开（设 $\boldsymbol{\eta}^R=\hat{\mathbf{X}}\mathbf{X}^{-1}\approx\mathrm{Exp}(\boldsymbol{\xi})$，创新量 $\mathbf{z}_n=\hat{\mathbf{X}}_n\mathbf{Y}_n-\mathbf{d}$）：
\begin{equation}
  \mathbf{z}_n=\hat{\mathbf{X}}_n\mathbf{Y}_n-\mathbf{d}\approx\mathbf{H}\boldsymbol{\xi}+\tilde{\mathbf{V}}_n,
  \label{eq:inekf-innov}
\end{equation}
对单个 landmark $\boldsymbol{\ell}_i$（world 坐标，$\mathrm{SE}_2(3)$ 嵌入齐次 $[\boldsymbol{\ell}_i;0;1]$），在\emph{源文献序} $\boldsymbol{\xi}=[\delta\boldsymbol{\theta};\delta\mathbf{v};\delta\mathbf{p}]$ 下，观测雅可比
\begin{equation}
  \mathbf{H}_i=\begin{bmatrix}-\boldsymbol{\ell}_i^\wedge&\mathbf{0}_{3\times3}&\mathbf{I}_3\end{bmatrix}
  \label{eq:inekf-H}
\end{equation}
\emph{只依赖已知参考向量 $\boldsymbol{\ell}_i$，完全不依赖估计与真值}\cite{barrau2017invariant}。这是 InEKF 一致性的“最后一块拼图”：传播 $\mathbf{A}$ 与观测 $\mathbf{H}$ \emph{都}独立于估计。

\begin{note}
\textbf{一处工程边界.}\;
“$\mathbf{H}$ 独立于估计”要配合“噪声协方差也固定/不变”才使增益完全独立于估计。若观测噪声 $\mathbf{N}_n$ 需经 $\mathrm{Ad}_{\hat{\mathbf{X}}}$ 把传感器坐标系噪声变换到滤波器坐标系，增益仍\emph{间接}温和依赖估计（\cref{sec:stability} 的稳定性 Lemma 量化此微扰不破坏稳定）。不同传感器的 $\mathbf{N}_n$ 原始坐标不同，\emph{不应}机械套用同一个 $\mathbf{R}^\top\mathbf{N}\mathbf{R}$ 公式。
\end{note}

\subsection{完整离散 RIEKF 算法}
\label{subsec:riekf-algo}

\begin{algo}[右不变 InEKF（$\mathrm{SE}_2(3)$，IMU $+$ 右不变观测）]
\label{algo:riekf}
\textbf{输入}：初始 $\hat{\mathbf{X}}_0,\mathbf{P}_0$；IMU 流 $\{\boldsymbol{\omega}_k,\mathbf{a}_k\}$（间隔 $\Delta t$）；右不变量测流 $\{\mathbf{Y}_n,\mathbf{d}_n\}$。\par
\medskip
\textbf{预测步（每个 IMU tick）：}
\begin{enumerate}
  \item \emph{状态前推}（群上 $\mathrm{Exp}$ 积分，保持 $\hat{\mathbf{R}}\in\SO(3)$ 无需重正交化）：
  \[
    \hat{\mathbf{R}}\leftarrow\hat{\mathbf{R}}\,\mathrm{Exp}(\boldsymbol{\omega}_k\Delta t),\quad
    \hat{\mathbf{v}}\leftarrow\hat{\mathbf{v}}+(\mathbf{g}+\hat{\mathbf{R}}\mathbf{a}_k)\Delta t,\quad
    \hat{\mathbf{p}}\leftarrow\hat{\mathbf{p}}+\hat{\mathbf{v}}\Delta t+\tfrac12(\mathbf{g}+\hat{\mathbf{R}}\mathbf{a}_k)\Delta t^2.
  \]
  \item \emph{误差转移}（仅依赖输入，幂零结构精确）：$\boldsymbol{\Phi}\leftarrow\mathbf{I}+\mathbf{A}\Delta t+\tfrac12\mathbf{A}^2\Delta t^2$，$\mathbf{A}$ 见 \cref{eq:se23-A}。
  \item \emph{协方差传播}：$\mathbf{P}\leftarrow\boldsymbol{\Phi}\mathbf{P}\boldsymbol{\Phi}^\top+\bar{\mathbf{Q}}_k$，其中 $\bar{\mathbf{Q}}_k=\int_0^{\Delta t}\!\exp(\mathbf{A}(\Delta t-\tau))\,\mathbf{G}_k\mathbf{Q}_c\mathbf{G}_k^\top\exp(\mathbf{A}(\Delta t-\tau))^\top d\tau$（可 Van Loan 精确离散；端点近似为 $\boldsymbol{\Phi}\mathbf{G}_k\mathbf{Q}_c\mathbf{G}_k^\top\boldsymbol{\Phi}^\top\Delta t$）。$\mathbf{G}_k$ 是 $6$ 维 IMU 白噪声到 $9$ 维误差切空间的嵌入雅可比（RIEKF 常含 $\mathrm{Ad}_{\hat{\mathbf{X}}}$）。
\end{enumerate}
\textbf{更新步（收到右不变量测 $\mathbf{Y}_n=\hat{\mathbf{X}}^{-1}\mathbf{d}_n+\mathbf{V}_n$ 时）：}
\begin{enumerate}
  \item \emph{不变创新量}：$\mathbf{z}\leftarrow\hat{\mathbf{X}}\mathbf{Y}_n-\mathbf{d}_n$（仅依赖 $\boldsymbol{\eta}^R$，\cref{eq:inekf-innov}）。
  \item \emph{增益}：$\mathbf{S}\leftarrow\mathbf{H}\mathbf{P}\mathbf{H}^\top+\hat{\mathbf{N}}_n$，$\mathbf{K}\leftarrow\mathbf{P}\mathbf{H}^\top\mathbf{S}^{-1}$，$\mathbf{H}$ 见 \cref{eq:inekf-H}（不依赖估计）。
  \item \emph{状态更新（群乘法，RIEKF \textbf{左乘}）}：$\hat{\mathbf{X}}\leftarrow\mathrm{Exp}\big((\mathbf{K}\mathbf{z})^\wedge\big)\,\hat{\mathbf{X}}$。
  \item \emph{协方差更新（Joseph 形式，保正定）}：$\mathbf{P}\leftarrow(\mathbf{I}-\mathbf{K}\mathbf{H})\mathbf{P}(\mathbf{I}-\mathbf{K}\mathbf{H})^\top+\mathbf{K}\hat{\mathbf{N}}_n\mathbf{K}^\top$。
\end{enumerate}
\end{algo}

\begin{derivation}[离散过程噪声 $\bar{\mathbf{Q}}_k$ 的三阶级数与正定性]
\label{deriv:Q-series}
\cref{algo:riekf} 第 3 步把过程噪声协方差写成积分形 $\bar{\mathbf{Q}}_k=\int_0^{\Delta t}\boldsymbol{\Phi}(\Delta t,\tau)\,\boldsymbol{\Sigma}\,\boldsymbol{\Phi}(\Delta t,\tau)^\top d\tau$，其中 $\boldsymbol{\Sigma}:=\mathbf{G}_k\mathbf{Q}_c\mathbf{G}_k^\top$ 为白噪声经嵌入雅可比映到切空间的功率谱密度。除 Van Loan 精确离散外，小 $\Delta t$ 下还可显式级数化——这在批量估计里反复复用、值得逐项写清\cite{barfoot2024state}。

\textbf{Step 1（按 $s$ 幂展开被积函数）.}\; 记 $s:=\Delta t-\tau$，由 \cref{eq:se23-A} 的幂零结构 $\boldsymbol{\Phi}(\Delta t,\tau)=\exp(\mathbf{A}s)=\mathbf{I}+\mathbf{A}s+\tfrac12\mathbf{A}^2s^2$（$\mathrm{SE}_2(3)$ 情形\emph{精确}），故
\[
  \boldsymbol{\Phi}\boldsymbol{\Sigma}\boldsymbol{\Phi}^\top
  =\boldsymbol{\Sigma}+s\big(\mathbf{A}\boldsymbol{\Sigma}+\boldsymbol{\Sigma}\mathbf{A}^\top\big)
  +s^2\big(\mathbf{A}\boldsymbol{\Sigma}\mathbf{A}^\top+\tfrac12\mathbf{A}^2\boldsymbol{\Sigma}+\tfrac12\boldsymbol{\Sigma}(\mathbf{A}^\top)^2\big)+\mathcal{O}(s^3).
\]

\textbf{Step 2（逐项积分 $\int_0^{\Delta t}s^n\,ds=\Delta t^{n+1}/(n+1)$）.}\;
\begin{equation}
  \bar{\mathbf{Q}}_k\approx\Delta t\,\boldsymbol{\Sigma}
  +\tfrac12\Delta t^2\big(\mathbf{A}\boldsymbol{\Sigma}+\boldsymbol{\Sigma}\mathbf{A}^\top\big)
  +\tfrac16\Delta t^3\big(2\mathbf{A}\boldsymbol{\Sigma}\mathbf{A}^\top+\mathbf{A}^2\boldsymbol{\Sigma}+\boldsymbol{\Sigma}(\mathbf{A}^\top)^2\big),
  \label{eq:Q-series}
\end{equation}
精确到 $\mathcal{O}(\Delta t^3)$。一阶项 $\Delta t\,\boldsymbol{\Sigma}$ 即端点近似；二阶、三阶项是 $\mathbf{A}$ 在区间内“搬运”噪声的修正。

\textbf{Step 3（保正定的实务要点）.}\; Barfoot 特别提示：数值上\emph{需保留到三阶}方能让 $\bar{\mathbf{Q}}_k$ 保持正定\cite{barfoot2024state}——只留一阶端点项在大 $\Delta t$ 或强机动下可能丢失正定、令 Cholesky 失败。最后，一条供批量估计复用的关系（把“预测误差”二阶矩等同于离散过程噪声）为
\begin{equation}
  E\big[(\boldsymbol{\Phi}\,\delta\mathbf{x}_{k-1}-\delta\mathbf{x}_k)(\boldsymbol{\Phi}\,\delta\mathbf{x}_{k-1}-\delta\mathbf{x}_k)^\top\big]=\bar{\mathbf{Q}}_k,
  \label{eq:Q-batch}
\end{equation}
它说明：在 \cref{sec:preint-se23} 把同一运动模型用于因子图时，预积分窗内逐小步复合出的 $\bar{\mathbf{Q}}_k$ 正是运动先验因子的代价权重 $\mathbf{Q}_k$，滤波与批量\emph{共享同一套}噪声机械。
\end{derivation}

\begin{pitfall}[更新步用错乘法方向]
RIEKF 误差定义为 $\boldsymbol{\eta}^R=\hat{\mathbf{X}}\mathbf{X}^{-1}$，更新必须在\emph{左}侧施加修正（\cref{algo:riekf} 第 3 步）才能正确减小右不变误差；LIEKF（$\boldsymbol{\eta}^L=\mathbf{X}^{-1}\hat{\mathbf{X}}$）则\emph{右}乘 $\hat{\mathbf{X}}\leftarrow\hat{\mathbf{X}}\,\mathrm{Exp}((\mathbf{K}\mathbf{z})^\wedge)$。记忆法则：\emph{误差“从哪边定义”，更新就“从哪边乘”}。写反的症状：前几步“看似正常”但逐渐发散、NEES 持续偏高。另：创新量符号与乘法方向必须配套——若 $\mathbf{z}\approx\mathbf{H}\boldsymbol{\xi}$ 估的是 $\boldsymbol{\eta}^R$ 的误差方向，修正要施加\emph{反向}（部分文献写 $\hat{\mathbf{X}}\leftarrow\mathrm{Exp}((-\mathbf{K}\mathbf{z})^\wedge)\hat{\mathbf{X}}$，符号被吸收进 residual 定义）。
\end{pitfall}

\begin{pitfall}[创新量里误用真值 $\mathbf{X}$]
实现只能用 $\hat{\mathbf{X}}$ 与量测 $\mathbf{Y}$ 构造残差。若写成 $\mathbf{X}^{-1}\mathbf{Y}-\mathbf{d}$（用真值），数学上会\emph{消掉}误差、滤波器“看似完美”实则无效。
\end{pitfall}

下面用最小 C++ 片段演示 $\mathrm{SE}_2(3)$ landmark 观测雅可比 \cref{eq:inekf-H} 的组装（仅 Jacobian 块，非完整滤波器）：

\begin{codebox}[C++]{SE\_2(3) landmark 观测雅可比组装（误差序 [theta, v, p]）}
#include <Eigen/Dense>

// skew(w) * v = w.cross(v)
Eigen::Matrix3d skew(const Eigen::Vector3d& w) {
    Eigen::Matrix3d W;
    W <<   0.0, -w.z(),  w.y(),
         w.z(),    0.0, -w.x(),
        -w.y(),  w.x(),    0.0;
    return W;
}

// Right-invariant landmark observation Jacobian H_i = [-skew(l_i), 0, I]
// l_world: landmark in world frame (already known reference vector)
Eigen::Matrix<double, 3, 9> riekfLandmarkJacobian(const Eigen::Vector3d& l_world) {
    Eigen::Matrix<double, 3, 9> H = Eigen::Matrix<double, 3, 9>::Zero();
    H.block<3, 3>(0, 0) = -skew(l_world);   // attitude block
    // velocity block (cols 3..5) stays zero for static landmark
    H.block<3, 3>(0, 6) = Eigen::Matrix3d::Identity();  // position block
    return H;   // independent of the state estimate -- the InEKF key property
}
\end{codebox}

\begin{practice}
\begin{enumerate}
  \item 在平移群 $\mathbb{R}$ 上构造最小反例：分别用创新量 $\mathbf{z}=\hat{x}-y$ 与 $\mathbf{z}=y-\hat{x}$，推导更新校正的正确符号，体会“residual $+$ 乘向 $+$ 误差定义”三者必须闭合。
  \item 把 \cref{algo:riekf} 改为 LIEKF（左不变观测 GPS $\mathbf{Y}=\hat{\mathbf{X}}\mathbf{e}$，$\mathbf{H}=[\mathbf{0},\mathbf{0},\mathbf{I}]$），写出创新量与右乘更新，指出哪几步符号变。
  \item 从 $\mathbf{y}=\mathbf{R}^\top(\boldsymbol{\ell}-\mathbf{p})$ 出发，对 $\delta\boldsymbol{\theta},\delta\mathbf{p}$ 做一阶展开，验证 \cref{eq:inekf-H} 的姿态块 $-\boldsymbol{\ell}^\wedge$ 与位置块 $\mathbf{I}$ 符号。
\end{enumerate}
\end{practice}

\subsection{另一条路：把标准 \texorpdfstring{$\SE(3)$}{SE(3)} EKF 代数变形为不变形}
\label{subsec:c4-bridge}

\paragraph{动机：不变 EKF 也可以“事后”从普通 EKF 推出。}
\cref{algo:riekf} 是\emph{原生}（native）RIEKF——直接用不变误差与群乘法更新写就。Barfoot 给出一条互补视角\cite{barfoot2024state}：取一个\emph{已经}写在 $\SE(3)$（左扰动）上的标准 EKF（其均值预测 $\check{\mathbf{T}}_k=\boldsymbol{\Xi}_k\hat{\mathbf{T}}_{k-1}$、观测 $\mathbf{y}_k=\mathbf{T}_k\mathbf{p}+\mathbf{n}_k$），\emph{纯代数}地把它变形成不变 EKF 的标准形。这条桥让人看清“不变 EKF 不是另起炉灶，而是同一滤波器换一个协方差存储坐标系”。

\paragraph{预测步：协方差改存“静止系”得 $\mathbf{F}'=\mathbf{1}$。}
标准式的运动雅可比 $\mathbf{F}_{k-1}=\mathrm{Ad}(\boldsymbol{\Xi}_k)=\check{\boldsymbol{\mathcal{T}}}_k\hat{\boldsymbol{\mathcal{T}}}_{k-1}^{-1}$（$\boldsymbol{\mathcal{T}}:=\mathrm{Ad}(\mathbf{T})$）虽已不依赖状态，仍是个非平凡变换。把协方差用伴随\emph{搬到静止（惯性）系}存储，定义 $\check{\mathbf{P}}'_k:=\check{\boldsymbol{\mathcal{T}}}_k^{-1}\check{\mathbf{P}}_k\check{\boldsymbol{\mathcal{T}}}_k^{-\top}$、$\mathbf{Q}'_k:=\check{\boldsymbol{\mathcal{T}}}_k^{-1}\mathbf{Q}_k\check{\boldsymbol{\mathcal{T}}}_k^{-\top}$，则协方差预测 $\check{\mathbf{P}}_k=\mathbf{F}_{k-1}\hat{\mathbf{P}}_{k-1}\mathbf{F}_{k-1}^\top+\mathbf{Q}_k$ 化为
\begin{equation}
  \check{\mathbf{P}}'_k=\mathbf{F}'_{k-1}\,\hat{\mathbf{P}}'_{k-1}\,\mathbf{F}'^{\top}_{k-1}+\mathbf{Q}'_k,\qquad \mathbf{F}'_{k-1}=\mathbf{1}.
  \label{eq:c4-Fprime}
\end{equation}
代价是估计依赖被\emph{搬进}了 $\mathbf{Q}'_k$ 的定义——但对“让雅可比变常数”这一目标可接受。

\paragraph{更新步：左不变创新量与伪测量雅可比 $\mathbf{G}'=\mathbf{p}^\odot$。}
关键恒等式 $(\check{\mathbf{T}}_k\mathbf{p})^\odot=\check{\mathbf{T}}_k\,\mathbf{p}^\odot\,\check{\boldsymbol{\mathcal{T}}}_k^{-1}$（即 $\odot$ 算子的伴随搬运律，见 \cref{ch:lie}）把原观测雅可比 $\mathbf{G}_k=(\check{\mathbf{T}}_k\mathbf{p})^\odot$ 的状态依赖\emph{剥离}到两侧伴随里。均值校正随之写成
\begin{equation}
  \hat{\mathbf{T}}_k=\check{\mathbf{T}}_k\,\exp\!\Big(\big(\mathbf{K}'_k\,\underbrace{\check{\mathbf{T}}_k^{-1}(\mathbf{y}_k-\check{\mathbf{y}}_k)}_{\text{左不变创新量}}\big)^\wedge\Big),
  \label{eq:c4-update}
\end{equation}
其中 $\check{\mathbf{T}}_k^{-1}(\mathbf{y}_k-\check{\mathbf{y}}_k)$ 称\emph{左不变创新量}，可视作一个替换原 $\mathbf{y}_k$ 的\textbf{伪测量}模型，其雅可比恰为常数 $\mathbf{G}'_k=\mathbf{p}^\odot$（\emph{不依赖估计}）。增益与协方差更新于是回到标准形 $\mathbf{K}'_k=\check{\mathbf{P}}'_k\mathbf{G}'^\top_k(\mathbf{G}'_k\check{\mathbf{P}}'_k\mathbf{G}'^\top_k+\mathbf{R}'_k)^{-1}$、$\hat{\mathbf{P}}'_k=(\mathbf{1}-\mathbf{K}'_k\mathbf{G}'_k)\check{\mathbf{P}}'_k$。

\begin{note}
\textbf{那“一处近似”及如何避免.}\;
协方差更新由 $\hat{\mathbf{P}}_k=\check{\boldsymbol{\mathcal{T}}}_k(\mathbf{1}-\mathbf{K}'_k\mathbf{G}'_k)\check{\boldsymbol{\mathcal{T}}}_k^{-1}\check{\mathbf{P}}_k$ 化到标准形 $\hat{\mathbf{P}}'_k=(\mathbf{1}-\mathbf{K}'_k\mathbf{G}'_k)\check{\mathbf{P}}'_k$ 时，需\emph{近似} $\check{\boldsymbol{\mathcal{T}}}_k^{-1}\hat{\mathbf{P}}_k\check{\boldsymbol{\mathcal{T}}}_k^{-\top}\approx\hat{\mathbf{P}}'_k$（更新前后伴随被当作同一个）——这是整条变形里\emph{唯一}的近似\cite{barfoot2024state}。Barfoot 指出：若一开始就用\emph{右}扰动（对应左不变误差）配左不变测量模型重新推导，可\emph{完全避开}这处近似、直达标准形。这正是 \cref{algo:riekf} 选右不变误差的深层理由：原生 RIEKF 天然无此近似，本小节的“代数变形”只是把它\emph{从普通 EKF 反推}出来的对照路径。\textbf{二者结论一致}（$\mathbf{F}'=\mathbf{1}$、$\mathbf{G}'$ 常数），不必同时实现。
\end{note}

% ════════════════════════════════════════════════════════════════
\section{一致性：InEKF 为何结构性保持不可观方向}
\label{sec:consistency-inekf}

\paragraph{动机：把“yaw 不该被看见”说清楚。}
\cref{sec:why-manifold} 的失败三（虚假可观测性）是 InEKF 最大的卖点所在。本节用可观测性语言（\cref{ch:slam_est} 的概率视角 $+$ \cref{ch:eskf} 的 \cref{thm:bias-overconfident}）说明 InEKF 如何\emph{结构性}地（而非靠调参）保持 yaw/平移不可观子空间。

\paragraph{不可观方向。}VIO/SLAM 中世界系无外部绝对参考时，以下变换不改变相机图像与 IMU 相对测量：(1) 所有位姿与地图点整体平移同一 $\mathbf{t}_0$；(2) 整体绕重力方向旋转同一 yaw 角 $\psi$。故理想估计问题至少有 $4$ 个规范（gauge）自由度（$3$ 平移 $+1$ yaw）。滤波器若声称这 $4$ 方向不确定性变小，就是利用了\emph{不存在}的信息。

用线性系统语言：误差系统 $\delta\mathbf{x}_{k+1}=\mathbf{F}_k\delta\mathbf{x}_k$、残差 $\mathbf{r}_k=\mathbf{H}_k\delta\mathbf{x}_k+\mathbf{n}_k$ 的不可观子空间 $\mathbf{N}$（其列张成不可观方向）应满足
\begin{equation}
  \mathbf{H}_k\mathbf{F}_{k-1}\cdots\mathbf{F}_0\,\mathbf{N}=\mathbf{0},\qquad\forall k.
  \label{eq:obs-null}
\end{equation}

\paragraph{标准 EKF 为何破坏它。}标准 EKF 在不断变化的估计点 $\hat{\mathbf{x}}_k$ 求 $\mathbf{F}_k,\mathbf{H}_k$。第 $i$ 帧在 $\hat{\mathbf{x}}_i$、第 $j$ 帧在 $\hat{\mathbf{x}}_j$ 线性化，每个局部雅可比单独看都合理，但堆叠起来\emph{不再共享同一 gauge 变换}，于是 $\mathbf{H}_k(\hat{\mathbf{x}}_k)\mathbf{F}_{k-1}(\hat{\mathbf{x}}_{k-1})\cdots\mathbf{N}\neq\mathbf{0}$——线性化系统\emph{误以为} yaw 影响残差。Kalman 更新 $\mathbf{P}^+=(\mathbf{I}-\mathbf{K}\mathbf{H})\mathbf{P}(\mathbf{I}-\mathbf{K}\mathbf{H})^\top+\mathbf{K}\mathbf{R}\mathbf{K}^\top$ 会沿 yaw 减小协方差，$\mathbf{N}^\top\mathbf{P}\mathbf{N}$ 下降——真实世界没提供 yaw 绝对信息、协方差却下降了，这就是一致性破坏（与 \cref{thm:bias-overconfident} 同源：线性化制造了虚假信息）。

\paragraph{InEKF 为何天然保持它。}对群仿射系统，误差对数满足 $\dot{\boldsymbol{\xi}}=\mathbf{A}_t\boldsymbol{\xi}$，$\mathbf{A}_t$ \emph{不依赖} $\hat{\mathbf{X}}_t$（\cref{eq:se23-A}）；不变/fundamental 观测的 $\mathbf{H}$ 也只依赖已知参考量（\cref{eq:inekf-H}）。于是对任何 gauge 方向 $\mathbf{N}$，只要物理观测满足 $\mathbf{H}\mathbf{N}=\mathbf{0}$，传播后仍有 $\mathbf{H}\boldsymbol{\Phi}(t,0)\mathbf{N}=\mathbf{0}$（因 $\boldsymbol{\Phi}$ 由仅依赖输入的 $\mathbf{A}_t$ 决定，不随估计漂移）。\emph{InEKF 保持一致性不靠调参、不靠小心选线性化点，而靠误差坐标与系统对称性一致}\cite{barrau2017invariant,hartley2020contact}——这正是 \cref{sec:lie-symmetry} 对称性地基的兑现。

\begin{insight}[一致性的真正根源]
一致性问题的根源\emph{不是} EKF 的一阶近似本身，而是“一阶近似\emph{破坏了原问题的对称性}”。InEKF 的价值在于让线性化后的误差系统仍然“记得”原问题有哪些方向不可观。这也解释了 FEJ：FEJ 把雅可比冻结在首次估计点，让所有残差\emph{被迫}共享同一 gauge 坐标——是“补丁式”地恢复对称；InEKF 则从误差定义上\emph{原生}尊重对称。
\end{insight}

\begin{pitfall}[FEJ 与 InEKF 叠加使用]
FEJ/OC-EKF 与 InEKF 都在解决同一个不一致问题。\emph{二选一}——同时使用两套一致性修复会重复修正，可能使协方差不升反降。InEKF 应\emph{单独}用（\cref{ch:eskf} 的 \cref{sec:filter-vs-opt} 提到的工程取舍）。
\end{pitfall}

\begin{practice}
\begin{enumerate}
  \item 对 landmark residual $\mathbf{r}=\mathbf{R}^\top(\boldsymbol{\ell}-\mathbf{p})-\mathbf{y}$，构造全局平移 gauge 方向 $\mathbf{N}_t$（$\mathbf{p}\mapsto\mathbf{p}+\boldsymbol{\epsilon}\mathbf{t}$、$\boldsymbol{\ell}\mapsto\boldsymbol{\ell}+\boldsymbol{\epsilon}\mathbf{t}$），用数值差分验证 $\mathbf{H}\mathbf{N}_t\approx\mathbf{0}$。
  \item 解释为什么 \cref{eq:obs-null} 在标准 EKF 下失败、在 InEKF 下保持（关键：$\mathbf{A}_t,\mathbf{H}$ 是否依赖 $\hat{\mathbf{X}}$）。
  \item （概念）若给系统加磁力计（提供绝对航向），yaw 还不可观吗？此时 InEKF 的一致性优势是否消失？
\end{enumerate}
\end{practice}

% ════════════════════════════════════════════════════════════════
\section{零偏的边界：群仿射 negative result 与修复路径}
\label{sec:bias-negative}

\paragraph{动机：完美 InEKF 有边界，IMU 零偏就在边界之外。}
\cref{ch:eskf} 的 \cref{paper:inekf} 已提示：把零偏增广进状态会\emph{破坏}群仿射，故工程用“不完美 InEKF”。本节给这个 negative result 的\emph{代数证明}与三条修复路径。

\begin{theorem}[零偏破坏群仿射（Barrau 2015 / Hartley 2020）]\label{thm:bias-negative}
在 $G=\mathrm{SE}_2(3)\times\mathbb{R}^6$（或其它 pose$\times$bias 直积/半直积）上，\emph{不存在}自然的 $f$ 使群仿射 \cref{eq:group-affine} 对所有元素成立。
\end{theorem}

\begin{derivation}[最小情形 $\SO(3)\times\mathbb{R}^3$（陀螺零偏）的代数证明]
\label{deriv:bias-negative}
取 $\mathbf{X}=(\mathbf{R},\mathbf{b}_g)$，群律 $(\mathbf{R}_1,\mathbf{b}_1)(\mathbf{R}_2,\mathbf{b}_2)=(\mathbf{R}_1\mathbf{R}_2,\mathbf{b}_1+\mathbf{b}_2)$，$\mathbf{e}=(\mathbf{I},\mathbf{0})$，动力学 $f_{\mathbf{u}}(\mathbf{R},\mathbf{b}_g)=(\mathbf{R}(\boldsymbol{\omega}_m-\mathbf{b}_g)^\wedge,\mathbf{0})$。代入 \cref{eq:group-affine}：
\begin{align*}
  \text{LHS}_{(1,1)}&=\mathbf{R}_1\mathbf{R}_2(\boldsymbol{\omega}_m-\mathbf{b}_1-\mathbf{b}_2)^\wedge,\\
  \text{RHS}_{(1,1)}&=\mathbf{R}_1(\boldsymbol{\omega}_m-\mathbf{b}_1)^\wedge\mathbf{R}_2+\mathbf{R}_1\mathbf{R}_2(\boldsymbol{\omega}_m-\mathbf{b}_2)^\wedge-\mathbf{R}_1\boldsymbol{\omega}_m^\wedge\mathbf{R}_2.
\end{align*}
用 $\boldsymbol{\omega}^\wedge\mathbf{R}=\mathbf{R}(\mathbf{R}^\top\boldsymbol{\omega})^\wedge$ 化简前两项与第三项之差：$\mathbf{R}_1(\boldsymbol{\omega}_m-\mathbf{b}_1)^\wedge\mathbf{R}_2-\mathbf{R}_1\boldsymbol{\omega}_m^\wedge\mathbf{R}_2=-\mathbf{R}_1\mathbf{b}_1^\wedge\mathbf{R}_2=-\mathbf{R}_1\mathbf{R}_2(\mathbf{R}_2^\top\mathbf{b}_1)^\wedge$，故 $\text{RHS}=\mathbf{R}_1\mathbf{R}_2\big[(\boldsymbol{\omega}_m-\mathbf{b}_2)^\wedge-(\mathbf{R}_2^\top\mathbf{b}_1)^\wedge\big]$。等式 $\text{LHS}=\text{RHS}$ 要求
\begin{equation}
  \mathbf{b}_1^\wedge=(\mathbf{R}_2^\top\mathbf{b}_1)^\wedge\quad\forall\mathbf{R}_2\in\SO(3),
  \label{eq:bias-contradiction}
\end{equation}
此需 $\mathbf{b}_1=\mathbf{0}$（取任意非单位 $\mathbf{R}_2$ 与 $\mathbf{b}_1\neq\mathbf{0}$ 即矛盾）。$\square$
\end{derivation}

\noindent{}\textbf{根因}：$\boldsymbol{\omega}\to\boldsymbol{\omega}-\mathbf{b}_g$ 引入“旋转与加性零偏\emph{不可交换}的耦合”，使群仿射在 $\mathbf{X}=\mathbf{e}$ 邻域即失败；加速度计零偏同理。

\paragraph{修复路径一：不完美 InEKF（imperfect InEKF）。}Hartley 等\cite{hartley2020contact} 的工程妥协：pose 部分用右不变误差 $\boldsymbol{\eta}=\hat{\mathbf{X}}\mathbf{X}^{-1}\in\mathrm{SE}_2(3)$、零偏用加性误差 $\tilde{\mathbf{b}}=\hat{\mathbf{b}}-\mathbf{b}\in\mathbb{R}^6$，联合线性化 $\boldsymbol{\epsilon}=(\boldsymbol{\xi},\tilde{\mathbf{b}})\in\mathbb{R}^{15}$。连续误差雅可比
\begin{equation}
  \mathbf{F}=\begin{bmatrix}\mathbf{F}_{\mathrm{pp}}(\mathbf{u})&\mathbf{F}_{\mathrm{pb}}(\hat{\mathbf{X}})\\\mathbf{0}&\mathbf{0}\end{bmatrix},
  \label{eq:imperfect-F}
\end{equation}
\emph{pose-pose 块} $\mathbf{F}_{\mathrm{pp}}$ 独立于估计（仍享 log-linear 优势），但\emph{pose-bias 块} $\mathbf{F}_{\mathrm{pb}}$ 形如 $\mathrm{diag}(-\hat{\mathbf{R}},-\hat{\mathbf{v}}^\wedge\hat{\mathbf{R}},-\hat{\mathbf{p}}^\wedge\hat{\mathbf{R}})$ \emph{依赖} $\hat{\mathbf{X}}$——这是群仿射被破坏的直接代数体现。代价：严格一致性不再，\cref{sec:stability} 的全局结论不直接成立；但 pose 主块的线性性仍带来巨大优势。

\paragraph{修复路径二、三：切群与等变滤波。}(二) \emph{切群/双参考系群}（Chauchat--Barrau--Bonnabel\cite{chauchat2022tangent}）在切群上重定义零偏群律（半直积），$\boldsymbol{\omega}=0$ 时恢复群仿射；(三) \emph{等变滤波 EqF}（\cref{sec:ukfm-ikfom}）完全绕开群仿射、只要等变性，把零偏通过对称群 lift 完全吸收——其 group action 让 $(\mathbf{R},\mathbf{b}_g)\mapsto\mathbf{R}^\top\mathbf{b}_g$ 方式耦合，恰好补偿 \cref{eq:bias-contradiction} 的不可交换性\cite{fornasier2022overcoming}。

\paragraph{Cassie 双足实证。}Hartley 等\cite{hartley2020contact} 在 Cassie 上做 $100$ 次 Monte-Carlo 仿真（相同噪声统计与初始协方差，但初始姿态/速度随机化——姿态在 $\pm30^\circ$、速度在 $\pm1.0\,$m/s 内均匀采样）：在相同接触信息下，右不变 InEKF 对\emph{所有可观状态}的收敛\emph{显著快于}四元数 EKF，收敛盆地（吸引域）明显更大。\footnote{原文（Hartley 等\cite{hartley2020contact}）给出的是\emph{定性}收敛对比（“可观状态收敛更快、吸引域更大”），并未报告位置/yaw RMSE 的具体百分比；且 IMU$+$接触下 yaw 本不可观（两滤波器的 yaw 均不收敛），故不存在“yaw 误差下降”这类稳态指标。需要精确数值请以原文图表为准。}性能差异的根因（\cref{sec:log-linear,sec:consistency-inekf}）：传播 $\mathbf{A}$ 与接触观测 $\mathbf{H}$ 都不依赖估计，零空间自动一致——\emph{同样的}接触信息在四元数 EKF 中无法享有这些结构性优势。

\begin{insight}[力量来自对称性匹配，而非“接触信息更多”]
接触辅助 InEKF 的力量\emph{不是}来自“接触信息更丰富”（四元数 EKF 也用同样接触信息），而是来自\emph{误差定义与群结构的匹配}：把接触点 $\mathbf{d}_i$ 嵌入 $\mathrm{SE}_{N+2}(3)$（每个接触点是“短寿命 landmark”，$\dot{\mathbf{d}}_i=\mathbf{0}$，与 SLAM 静态地标 $\dot{\boldsymbol{\ell}}_j=\mathbf{0}$ 数学结构一致）后，接触观测自然成为 fundamental 右不变观测，享有 \cref{sec:stability} 的全部保证。这把 IMU 导航、landmark SLAM、接触里程计压到\emph{同一个}代数模板。
\end{insight}

\begin{practice}
\begin{enumerate}
  \item 独立复现 \cref{deriv:bias-negative}，指出导致矛盾的关键步 \cref{eq:bias-contradiction}（$\mathbf{b}_1^\wedge\neq(\mathbf{R}_2^\top\mathbf{b}_1)^\wedge$）。
  \item 写出 $\mathrm{SE}_4(3)$（四足四接触点）的矩阵形式与群乘法，验证 $\dot{\mathbf{d}}_i=\mathbf{0}$ 对应的动力学仍群仿射。
  \item （概念）解释为什么“InEKF 的 $\mathbf{A}$ 不依赖估计 $\Rightarrow$ 带零偏的 InEKF 也一定全面优于 ESKF”是\emph{错}的（提示：$\mathbf{F}_{\mathrm{pb}}$ 块，pose 享优势但 bias 块退化为普通 EKF 水平）。
\end{enumerate}
\end{practice}

% ════════════════════════════════════════════════════════════════
\section{免雅可比与推广：UKF-M、IKFoM、EqF}
\label{sec:ukfm-ikfom}

\paragraph{动机：当雅可比难算、或状态不是李群时。}
InEKF 要求状态是李群、动力学群仿射。现实里两条件常不满足：雅可比难推（UKF-M 的用武之地）、状态在齐性空间（EqF 的用武之地）、或只想要通用混合流形（IKFoM）。本节把流形滤波族补齐。

\subsection{UKF-M：免雅可比的 sigma 点流形滤波}
\label{subsec:ukfm}

\paragraph{核心思想（Brossard--Barrau--Bonnabel 2020）。}在切空间（李代数）生成 sigma 点 $\to$ 经 retraction $\varphi$ 映到流形 $\to$ 传播/观测 $\to$ 经 $\varphi^{-1}$ 拉回切空间算均值/协方差\cite{brossard2020ukfm}。\emph{好处}：免解析雅可比；\emph{代价}：每步需 $2(d+q)+1$ 次 $f/h$ 求值（$d$ 状态维、$q$ 噪声维）。不确定度表示 $\mathbf{X}=\varphi(\hat{\mathbf{X}},\boldsymbol{\xi})$，$\boldsymbol{\xi}\sim\mathcal{N}(\mathbf 0,\mathbf{P})$，$\varphi$ 满足 $\varphi(\hat{\mathbf{X}},\mathbf 0)=\hat{\mathbf{X}}$ 且 $\partial_{\boldsymbol{\xi}}\varphi(\hat{\mathbf{X}},\mathbf 0)=\mathbf{I}$。右扰动 retraction $\varphi(\hat{\mathbf{X}},\boldsymbol{\xi})=\hat{\mathbf{X}}\,\mathrm{Exp}(\boldsymbol{\xi})$（对应左不变误差），左扰动对应右不变误差。

\begin{insight}[UKF-M 的两个关键简化]
(1) \emph{均值}不通过流形上的迭代 Fréchet 均值，而\emph{直接}用无噪声传播 $\hat{\mathbf{X}}^-=f(\hat{\mathbf{X}},\boldsymbol{\omega},\mathbf 0)$——这是 UKF-M 相较前人的核心工程贡献\cite{brossard2020ukfm}。(2) sigma 点从一开始就“住在流形上”（经 $\mathrm{Exp}$ 自然生成的合法群元素），不存在“非单位旋转/非归一四元数”问题；标准 UKF 在环境空间（如四元数 $\mathbb{R}^4$）生成 sigma 点须事后归一化，破坏 sigma 点的统计对称性。
\end{insight}

\begin{pitfall}[用欧氏平均代替 Fréchet 均值]
对 sigma 点的旋转部分直接算 $\bar{\mathbf{R}}=\sum w_j\mathbf{R}_j$（这\emph{不是}合法旋转矩阵！）或对角度直接算术平均，当 sigma 点跨越 $\pm\pi$ 时均值突然跳到 $0$、协方差特征值爆炸。根因：$\SO(2)$/$\SO(3)$ 不是向量空间，加法平均无几何意义。UKF-M 用无噪声传播取代迭代 Fréchet 均值；若需精确均值，对 sigma 点迭代 $\bar{\mathbf{X}}\leftarrow\bar{\mathbf{X}}\,\mathrm{Exp}\big(\sum w_j\mathrm{Log}(\bar{\mathbf{X}}^{-1}\mathbf{X}_j)\big)$（$3$--$5$ 次收敛）。
\end{pitfall}

\paragraph{与 InEKF 的关系。}UKF-M 免雅可比、适用任意可平行化流形（不需群仿射），但\emph{不享有} InEKF 的“误差动力学与状态解耦”，一致性保持不如 InEKF——\emph{两者互补而非替代}。在 $\mathrm{SE}_2(3)$ 惯导实验中，$\mathrm{SE}_2(3)$-UKF（右扰动）在 $45^\circ$ 初始航向误差下显著优于 $\SO(3)\times\mathbb{R}^6$-UKF 与原版 EKF\cite{brossard2020ukfm}。

\subsection{IKFoM 与 FAST-LIO2：iterated 不是 invariant}
\label{subsec:ikfom}

\begin{pitfall}[“IEKF” 的命名冲突——头号阅读陷阱]
看到 “IEKF” 务必先确认是哪种：\textbf{Iterated} EKF（Bell--Cathey 1993，同一观测更新内反复线性化，等价单步 Gauss--Newton，\cref{ch:eskf} 的 \cref{eq:iekf} 已给）与 \textbf{Invariant} EKF（Barrau--Bonnabel 2017，本章主题）\emph{完全不同}。\textbf{IKFoM} 的 “I” 是 \emph{iterated}（He--Xu--Zhang 2021\cite{he2021ikfom}），\emph{不是} invariant。二者在\emph{正交}维度改进经典 EKF：iterated 解决“测量非线性强、一次线性化不够”，invariant 解决“误差坐标不对、稳定性/可观测性被破坏”。
\end{pitfall}

\paragraph{IKFoM 核心。}对混合流形 $\mathcal{S}=\mathbb{R}^m\times\SO(3)\times\cdots\times S^2$，用户只需定义每个子流形的 $\boxplus/\boxminus$，系统自动合成 iterated ESKF\cite{he2021ikfom}：$\mathbf{x}_{k+1}=\mathbf{x}_k\boxplus(\Delta t\,f(\mathbf{x}_k,\mathbf{u}_k,\mathbf{w}_k))$，$\SO(3)$ 的 $\boxplus$ 默认右乘 $\mathbf{R}\boxplus\boldsymbol{\phi}=\mathbf{R}\,\mathrm{Exp}(\boldsymbol{\phi})$（与本书右扰动一致）。FAST-LIO2\cite{xu2022fastlio2} 的状态 $\SO(3)\times\mathbb{R}^{15}\times S^2\times\SO(3)\times\mathbb{R}^3$ 是“分离式”product manifold、与 $\mathrm{SE}_2(3)$ 矩阵群\emph{不等价}，其传播 $\mathbf{F}$ 依赖估计 $\hat{\mathbf{R}}$，\emph{不}满足 log-linear——故 FAST-LIO2 在理论上是一般 EKF-on-manifold 的局部收敛保证、\emph{不是} InEKF 的轨迹无关吸引域。但 LiDAR 信息丰富，这一理论差距在实际中通常不显著。其快增益公式 $\mathbf{K}=(\mathbf{H}^\top\mathbf{R}^{-1}\mathbf{H}+\mathbf{P}^{-1})^{-1}\mathbf{H}^\top\mathbf{R}^{-1}$ 维度为\emph{状态}（$\sim24$）而非测量（$\sim$ 数千），避免大矩阵求逆。

\subsection{（选读）等变滤波 EqF：InEKF 的严格推广}
\label{subsec:eqf}

\paragraph{从群仿射到等变。}EqF（Mahony--van Goor--Hamel\cite{mahony2021equivariant}，\cref{ch:lie} 的 \cref{sec:lie-symmetry} 已提）把 InEKF 推广到\emph{齐性空间} $\mathcal{M}=G/H$：状态 $\boldsymbol{\xi}\in\mathcal{M}$ 未必是李群，只需存在李群 $G$ 在 $\mathcal{M}$ 上的传递作用 $\phi$ 使动力学\emph{等变} $D\phi_{\mathbf{X}}\cdot f(\boldsymbol{\xi},\mathbf{u})=f(\phi_{\mathbf{X}}(\boldsymbol{\xi}),\psi_{\mathbf{X}}(\mathbf{u}))$（\cref{def:equivariance} 的动力学版）。

\begin{table}[H]\centering\small
\caption{InEKF vs EqF：条件与结论}
\label{tab:inekf-eqf}
\begin{tabular}{lll}
\toprule
属性 & InEKF (2017) & EqF (2022) \\
\midrule
状态空间 & 李群 $G$ & 齐性空间 $\mathcal{M}=G/H$ \\
代数条件 & group-affine & equivariance（更弱） \\
线性化点 & 沿估计轨迹（不变误差） & 固定原点 $\boldsymbol{\xi}^\circ$ \\
$\mathbf{F}$ 独立估计 & 传播段是 & 总是 \\
$\mathbf{H}$ 独立估计 & fundamental 观测时是 & 输出等变时是 \\
处理零偏 & 不完美（失真，\cref{eq:imperfect-F}） & 通过 lift 恢复一致（\cite{fornasier2022overcoming}） \\
\bottomrule
\end{tabular}
\end{table}

\begin{insight}[直觉类比]
InEKF 要求“状态空间本身是群”——像要求地球表面是平面；EqF 只要求“有一个群在状态空间上作用”——承认地球是球面、但用旋转群在球面上的作用简化问题。球面不是群（两点无天然“乘法”），但旋转群能把任一点移到任一点（传递性），这就够了。代价是要构造 lift 映射 $\Lambda:\mathcal{M}\times\mathbb{V}\to\mathfrak{g}$ 把动力学“抬升”到对称群。EqF 进一步利用输出等变性把输出误差分解为等变（精确）$+$ 非等变（高阶小量）部分，得\emph{三阶}线性化误差（标准 EKF 仅二阶）；EqVIO\cite{vangoor2023eqvio} 在 EuRoC 上 $5/11$ 序列取得最佳 RMSE、速度比第二快算法快约 $2$ 倍。
\end{insight}

\begin{practice}
\begin{enumerate}
  \item 写出 UKF-M 的 $\SO(3)$ product retraction $\varphi/\varphi^{-1}$（提示：$\varphi(\mathbf{R},\boldsymbol{\phi})=\mathbf{R}\,\mathrm{Exp}(\boldsymbol{\phi})$、$\varphi^{-1}(\hat{\mathbf{R}},\mathbf{R})=\mathrm{Log}(\hat{\mathbf{R}}^\top\mathbf{R})$），说明它为何天然避开角度缠绕。
  \item 用一句话解释为什么 FAST-LIO2 用 IKFoM 而非 InEKF（提示：状态含 $S^2$ 与双外参、product manifold 非 $\mathrm{SE}_2(3)$、不满足群仿射）。
  \item （概念）EqF 如何用 lift 处理零偏，从而绕开 \cref{thm:bias-negative} 的 negative result？
\end{enumerate}
\end{practice}

% ════════════════════════════════════════════════════════════════
\section{（选读）稳定性：收敛半径独立于初始化时刻}
\label{sec:stability}

\paragraph{动机：为什么 InEKF 的稳定性是“真定理”。}
log-linear（\cref{sec:log-linear}）保证传播精确，但更新步仍有非线性。Barrau--Bonnabel 的第三个贡献是证明 InEKF 是\emph{渐近稳定观测器}，且收敛半径\emph{独立于初始化时刻 $t_0$}——这是标准 EKF 稳定性结论永远达不到的。

\begin{theorem}[InEKF 渐近稳定（Barrau--Bonnabel TAC 2017, Thm 4）]\label{thm:inekf-stable}
考虑群仿射系统 \cref{eq:ga-dynamics} 配左（右）不变观测。若 Deyst--Price 1968 的稳定性条件（一步转移可逆、过程/观测噪声正定、过程/观测 Gramian 一致上下有界，共五条）沿\emph{真实轨迹 $\mathbf{X}_t$} 线性化后成立，则 LIEKF（RIEKF）的估计 $\hat{\mathbf{X}}_t$ 是 $\mathbf{X}_t$ 的渐近稳定观测器，且收敛半径 $\varepsilon>0$ 在整条轨迹上\emph{一致}（独立于初始化时间 $t_0$）。
\end{theorem}

\paragraph{证明思路（附录 C 的骨架）。}核心三步\cite{barrau2017invariant}：(1) 更新阶段 $\boldsymbol{\xi}^+_n=(\mathbf{I}-\mathbf{L}_n\mathbf{H})\boldsymbol{\xi}_n+\mathbf{r}_n(\boldsymbol{\xi}_n)$，由 BCH 知余项 $\mathbf{r}_n=\mathcal{O}(\|\boldsymbol{\xi}_n\|\|\mathbf{L}_n\boldsymbol{\xi}_n\|)$ 即\emph{二阶}；(2) Duhamel 分解 $\boldsymbol{\xi}_t=\boldsymbol{\Psi}^t_{t_0}\boldsymbol{\xi}_0+\sum_{t_n<t}\boldsymbol{\Psi}^t_{t_n}\mathbf{r}_n(\boldsymbol{\xi}_{t_n})$；(3) 由 Deyst--Price 的 Lyapunov 函数 $V=\boldsymbol{\xi}^\top\mathbf{P}^{-1}\boldsymbol{\xi}$ 得线性主项严格单减、二阶余项被一致界吸收。\emph{关键}：因 $\mathbf{A}_{\mathbf{u}_t}$ 只依赖输入，只要输入 $\mathbf{u}_t$ 整条轨迹\emph{一致有界}，BCH 高阶余项就能被\emph{单一常数}控制——$\varepsilon$ 只需小到此常数水平、与 $t_0$ 无关。

\begin{insight}[$\varepsilon\perp t_0$ 为何不平凡]
Reif--Unbehauen 1999 的 EKF 稳定性给出 $\varepsilon(t_0)\le\text{const}/(\sup_{t\ge t_0}\|\mathbf{A}(\hat{\mathbf{x}}_t,\mathbf{u}_t)\|\cdot\gamma_{\max})$，$\varepsilon$ \emph{随 $t_0$ 变化}——因真实轨迹上雅可比幅度随时间变、二阶余项随之膨胀。机器人运行很久后，$\hat{\mathbf{x}}_t$ 可能进入与初始证明区域\emph{完全不同}的状态区域，早期收敛半径不再适用。InEKF 传播段无二阶余项、余项只来自更新，Deyst--Price 型 LTV-KF 稳定性可在\emph{任意}时间窗口重复使用，故收敛半径不依赖“从何时开始滤波”。
\end{insight}

\begin{table}[H]\centering\small
\caption{三大定理的逻辑链（Barrau--Bonnabel TAC 2017）}
\label{tab:three-theorems}
\begin{tabular}{llll}
\toprule
环节 & 定理 & 从 & 到 \\
\midrule
代数 $\to$ 微分几何 & \cref{thm:ga-equiv}（群仿射） & 代数恒等式 \cref{eq:group-affine} & 误差 ODE 自治 \\
微分几何 $\to$ 分析 & \cref{thm:log-linear}（log-linear） & 流的群同态 \cref{eq:flow-homo} & 对数坐标精确线性 \\
分析 $\to$ 控制 & \cref{thm:inekf-stable}（稳定性） & Deyst--Price LTV-KF & $\varepsilon\perp t_0$ \\
\bottomrule
\end{tabular}
\end{table}

\begin{note}
\textbf{不是全局收敛.}\;
\cref{thm:inekf-stable} 仍是\emph{局部}收敛（半径 $\varepsilon$ 内），只是 $\varepsilon$ 独立 $t_0$。$\log$ 映射有单射半径限制、观测可能退化、接触误检会破坏模型——\emph{切勿}把它读成“InEKF 全局收敛”。其精确意义：在可观性 $+$ 局部误差条件满足时，收敛邻域不随初始化时刻漂移，这是普通 EKF 难以保证的。
\end{note}

\begin{practice}
\begin{enumerate}
  \item 写出 Deyst--Price 的 Lyapunov 函数 $V=\boldsymbol{\xi}^\top\mathbf{P}^{-1}\boldsymbol{\xi}$，对无噪声 log-linear 系统 $\dot{\boldsymbol{\xi}}=\mathbf{A}_t\boldsymbol{\xi}$ 求 $\dot V$，说明持续观测下 $V$ 单减。
  \item 对比 \cref{thm:inekf-stable} 与 Reif--Unbehauen 1999：雅可比依赖性、二阶余项、$\varepsilon$ 与 $t_0$ 关系三处差异。
  \item （概念）举例说明 $\varepsilon\perp t_0$ 的工程价值（提示：摔倒后重站、长时间巡检、接触频繁切换）。
\end{enumerate}
\end{practice}

% ════════════════════════════════════════════════════════════════
\section{\texorpdfstring{$\mathrm{SE}_2(3)$}{SE2(3)} 不变预积分：群仿射机械的批量估计收尾}
\label{sec:preint-se23}

\paragraph{动机：把滤波器里的群仿射机械，原封不动搬到因子图。}
前面各节把 $\mathrm{SE}_2(3)$ 群仿射结构用于\emph{滤波}（\cref{sec:inekf-algo}）。但 \cref{ch:eskf} 的 \cref{sec:filter-vs-opt} 已论证：批量/因子图（\cref{ch:nlopt}）往往精度更高。问题是——若在因子图里给\emph{每个} IMU 采样都放一个状态变量，变量数会爆炸。\textbf{预积分}（pre-integration）的思想正是把一窗 IMU 测量压成\emph{单一}运动先验因子，只在窗口两端 $\tau_0,\tau_J$ 保留状态\cite{forster2017preintegration}。\cref{ch:imu} 已给主流 Forster $\SO(3)$ 预积分；本节给其 $\mathrm{SE}_2(3)$ 变体（Brossard 等\cite{brossard2021associating}）——它把本章自建的 $\mathrm{SE}_2(3)$ 李机械（\cref{def:se23}）、群仿射（\cref{sec:group-affine}）、精确指数（\cref{eq:se23-Phi}）汇成一个\emph{天然融合}的收尾：\emph{同一套}扩展位姿群既驱动 InEKF、又驱动预积分因子。

\paragraph{反面：不用 $\mathrm{SE}_2(3)$ 会怎样。}
Forster 预积分把姿态放 $\SO(3)$、速度/位置放欧氏，三者分开传播、偏置修正用一阶链式。可行，但速度—位置—姿态的耦合被拆散，偏置重线性化（re-linearization）时要分别维护多块一阶修正。$\mathrm{SE}_2(3)$ 把三者并进\emph{一个}群元素，整窗增量成为群乘积 $\prod\boldsymbol{\Delta}_j\boldsymbol{\Xi}_j$，偏置雅可比经\emph{一个} $\mathrm{SE}_2(3)$ 左雅可比统一搬运——代数更紧、与本章 InEKF 共享零件。

\begin{note}
\textbf{本节排序约定（务必看清，与本章前文不同）.}\;
本节切空间向量统一采用 Barfoot/Brossard 的 $\mathrm{SE}_2(3)$ 顺序 $[\,\text{位置};\text{速度};\text{旋转}\,]$，\emph{与本书 \cref{def:se23} 的 $[\boldsymbol{\rho};\boldsymbol{\nu};\boldsymbol{\phi}]$ 一致}（位置在前、旋转在后）。这\emph{不同于} \cref{sec:log-linear} 的系统矩阵 \cref{eq:se23-A} 所用“旋转在前” $[\delta\boldsymbol{\theta};\delta\mathbf{v};\delta\mathbf{p}]$ 源文献序——后者是为与 Barrau--Bonnabel 的 $\mathbf{A}$ 矩阵逐块对照（见本章开头记号 note）。两序仅差置换矩阵 $\mathbf{P}$。群元素仍按 \cref{eq:imu-matrix-f}：$\mathbf{T}=\big[\begin{smallmatrix}\mathbf{R}&\mathbf{r}&\mathbf{v}\\\mathbf{0}^\top&1&0\\\mathbf{0}^\top&0&1\end{smallmatrix}\big]$（位置列 $\mathbf{r}$ 在前、速度列 $\mathbf{v}$ 在后）。
\end{note}

\subsection{精确均值传播与 \texorpdfstring{$\mathbf{N}$}{N} 矩阵}
\label{subsec:preint-mean}

\paragraph{有限区间的精确积分。}
在预积分窗内取局部时刻 $\tau_0,\dots,\tau_J$。设单步 $(\tau_{j-1},\tau_j)$ 内偏置与 IMU 读数恒定，记 $\bar{\boldsymbol{\omega}}_j=\tilde{\boldsymbol{\omega}}(\tau_j)-\mathbf{b}_\omega$、$\bar{\mathbf{a}}_j=\tilde{\mathbf{a}}(\tau_j)-\mathbf{b}_a$、$\Delta\tau_j=\tau_j-\tau_{j-1}$，则 \cref{eq:imu-kin} 退化为 LTI 系统、可\emph{精确}积分\cite{barfoot2024state}：
\begin{align}
  \bar{\mathbf{R}}(\tau_j)&=\bar{\mathbf{R}}(\tau_{j-1})\,\mathrm{Exp}\!\big(\Delta\tau_j\,\bar{\boldsymbol{\omega}}_j\big),\label{eq:preint-Rprop}\\
  \bar{\mathbf{v}}(\tau_j)&=\bar{\mathbf{v}}(\tau_{j-1})+\Delta\tau_j\big(\boldsymbol{\alpha}_v(\tau_j)+\mathbf{g}\big),\nonumber\\
  \bar{\mathbf{r}}(\tau_j)&=\bar{\mathbf{r}}(\tau_{j-1})+\Delta\tau_j\,\bar{\mathbf{v}}(\tau_{j-1})+\tfrac12\Delta\tau_j^2\big(\boldsymbol{\alpha}_r(\tau_j)+\mathbf{g}\big),\nonumber
\end{align}
其中比力的姿态搬运分两路——速度走 $\SO(3)$ \emph{左雅可比} $\mathbf{J}$、位置走\emph{新矩阵} $\mathbf{N}$：
\begin{equation}
  \boldsymbol{\alpha}_v(\tau_j)=\bar{\mathbf{R}}(\tau_{j-1})\,\mathbf{J}\!\big(\Delta\tau_j\,\bar{\boldsymbol{\omega}}_j\big)\,\bar{\mathbf{a}}_j,\qquad
  \boldsymbol{\alpha}_r(\tau_j)=\bar{\mathbf{R}}(\tau_{j-1})\,\mathbf{N}\!\big(\Delta\tau_j\,\bar{\boldsymbol{\omega}}_j\big)\,\bar{\mathbf{a}}_j.
  \label{eq:preint-alpha}
\end{equation}

\begin{definition}[$\mathbf{N}$ 矩阵：$\mathrm{SE}_2(3)$ 位置均值的“二次左雅可比”]\label{def:N-matrix}
对转角 $\phi$、单位轴 $\mathbf{a}$（$\boldsymbol{\phi}=\phi\mathbf{a}$），定义
\begin{equation}
  \mathbf{N}(\phi\mathbf{a})=2\!\int_0^1\!\alpha\,\mathbf{J}(\alpha\phi\mathbf{a})\,d\alpha
  =2\sum_{n=0}^\infty\frac{1}{(n+2)!}\big(\phi\mathbf{a}^\wedge\big)^n,
  \label{eq:N-series}
\end{equation}
其闭式为
\begin{equation}
  \mathbf{N}(\phi\mathbf{a})=2\,\frac{1-\cos\phi}{\phi^2}\,\mathbf{I}
  +\Big(1-2\,\frac{1-\cos\phi}{\phi^2}\Big)\mathbf{a}\mathbf{a}^\top
  +2\,\frac{\phi-\sin\phi}{\phi^2}\,\mathbf{a}^\wedge.
  \label{eq:N-closed}
\end{equation}
当 $\phi\to0$ 时 $\mathbf{N}\to\mathbf{I}$。它与 $\SO(3)$ 左雅可比 $\mathbf{J}(\phi\mathbf{a})=\int_0^1\mathrm{Exp}(\alpha\phi\mathbf{a})\,d\alpha=\sum_n\frac{1}{(n+1)!}(\phi\mathbf{a}^\wedge)^n$ 的区别仅多一个权重 $2\alpha$（级数里 $(n+1)!\to(n+2)!$）——这正对应位置是比力的\emph{二次}积分、速度是\emph{一次}积分。
\end{definition}

\begin{insight}[小步法为何“看不见” $\mathbf{N}$]
\cref{algo:riekf} 的状态前推与 \cref{ch:gins} 的 \cref{eq:gins-nominal} 都用 $\tfrac12(\mathbf{g}+\hat{\mathbf{R}}\mathbf{a})\Delta t^2$ 推位置——这\emph{隐含} $\mathbf{N}\approx\mathbf{J}\approx\mathbf{I}$，即 $\Delta\tau\to0$ 的极限。$\Delta\tau\,\bar{\boldsymbol{\omega}}$ 越大（高速旋转或大步长），$\mathbf{N},\mathbf{J}$ 偏离 $\mathbf{I}$ 越多，\cref{eq:preint-alpha} 才是\emph{无近似}的均值传播。换言之：\emph{主流小步积分是 $\mathrm{SE}_2(3)$ 精确积分在 $\mathbf{N}=\mathbf{J}=\mathbf{I}$ 处的一阶简化}——这解释了为何高频 IMU（$\Delta\tau$ 小）下小步法足够，而长预积分窗或低频 IMU 下应保留 $\mathbf{N},\mathbf{J}$。
\end{insight}

\subsection{（选读）时间机器：把均值传播写成 \texorpdfstring{$\mathrm{SE}_2(3)$}{SE2(3)} 三元积}
\label{subsec:preint-timemachine}

\paragraph{为什么需要一个“辅助群”。}
\cref{eq:preint-Rprop} 的位置更新含 $\Delta\tau_j\bar{\mathbf{v}}(\tau_{j-1})$ 这一\emph{跨步}耦合（当前位置依赖上一步速度），使整窗传播\emph{不能}直接写成纯 $\mathrm{SE}_2(3)$ 群乘积。Barfoot 引入一族辅助矩阵——\emph{时间机器}（time machines）——来吸收这个耦合\cite{barfoot2024state}。

\begin{definition}[时间机器群 $\mathbb{T}$]\label{def:time-machine}
$5\times5$ 矩阵 $\boldsymbol{\Delta}(\tau)=\big[\begin{smallmatrix}\mathbf{I}&\mathbf{0}&\mathbf{0}\\\mathbf{0}^\top&1&0\\\mathbf{0}^\top&\tau&1\end{smallmatrix}\big]$（$\tau\in\mathbb{R}$ 为时间增量）构成矩阵群 $\mathbb{T}$：含单位元（$\tau=0$）、逆 $\boldsymbol{\Delta}(\tau)^{-1}=\boldsymbol{\Delta}(-\tau)$、封闭 $\boldsymbol{\Delta}(\tau_1)\boldsymbol{\Delta}(\tau_2)=\boldsymbol{\Delta}(\tau_1+\tau_2)$。其 $9\times9$ 伴随版（作用于 $[\,\text{位置};\text{速度};\text{旋转}\,]$ 切向量）为
\begin{equation}
  \boldsymbol{\mathcal{D}}(\boldsymbol{\Delta})=\begin{bmatrix}\mathbf{I}&\tau\mathbf{I}&\mathbf{0}\\\mathbf{0}&\mathbf{I}&\mathbf{0}\\\mathbf{0}&\mathbf{0}&\mathbf{I}\end{bmatrix},
  \label{eq:time-machine-D}
\end{equation}
即把“位置切分量”加上 $\tau$ 倍“速度切分量”。
\end{definition}

\begin{derivation}[把单步均值传播写成 $\mathrm{SE}_2(3)$ 三元积]
\label{deriv:preint-triple}
受 Brossard 等\cite{brossard2021associating} 启发，\cref{eq:preint-Rprop} 的单步更新可\emph{精确}重写为
\begin{equation}
  \mathbf{T}_j=\mathbf{T}_{g,j}\,\boldsymbol{\Delta}_j^{-1}\mathbf{T}_{j-1}\boldsymbol{\Delta}_j\,\boldsymbol{\Xi}_j,
  \label{eq:preint-meanprop}
\end{equation}
其中三个 $\mathrm{SE}_2(3)$ 元素分别承载重力、上一步位姿（共轭后）、本步 IMU 增量：
\begin{equation}
  \mathbf{T}_{g,j}=\begin{bmatrix}\mathbf{I}&\tfrac12\Delta\tau_j^2\mathbf{g}&\Delta\tau_j\mathbf{g}\\\mathbf{0}^\top&1&0\\\mathbf{0}^\top&0&1\end{bmatrix},\quad
  \boldsymbol{\Xi}_j=\begin{bmatrix}\mathrm{Exp}(\Delta\tau_j\bar{\boldsymbol{\omega}}_j)&\tfrac12\Delta\tau_j^2\mathbf{N}_j\bar{\mathbf{a}}_j&\Delta\tau_j\mathbf{J}_j\bar{\mathbf{a}}_j\\\mathbf{0}^\top&1&0\\\mathbf{0}^\top&0&1\end{bmatrix},
  \label{eq:preint-Tg-Xi}
\end{equation}
$\mathbf{N}_j=\mathbf{N}(\Delta\tau_j\bar{\boldsymbol{\omega}}_j)$、$\mathbf{J}_j=\mathbf{J}(\Delta\tau_j\bar{\boldsymbol{\omega}}_j)$。\emph{时间机器的作用}：共轭 $\boldsymbol{\Delta}_j^{-1}\mathbf{T}_{j-1}\boldsymbol{\Delta}_j=\big[\begin{smallmatrix}\mathbf{R}_j&\mathbf{r}_j+\Delta\tau_j\mathbf{v}_j&\mathbf{v}_j\\\mathbf{0}^\top&1&0\\\mathbf{0}^\top&0&1\end{smallmatrix}\big]$ 恰好把跨步耦合 $\Delta\tau_j\mathbf{v}_j$ 注入位置列、且结果仍是 $\mathrm{SE}_2(3)$ 元素。$\boldsymbol{\Xi}_j$ 又可写成纯指数 $\boldsymbol{\Xi}_j=\mathrm{Exp}(\boldsymbol{\psi}_j)$，
\begin{equation}
  \boldsymbol{\psi}_j=\big[\,\tfrac12\Delta\tau_j^2\mathbf{J}_j^{-1}\mathbf{N}_j\bar{\mathbf{a}}_j\;;\;\;\Delta\tau_j\bar{\mathbf{a}}_j\;;\;\;\Delta\tau_j\bar{\boldsymbol{\omega}}_j\,\big],
  \label{eq:preint-psi}
\end{equation}
这一形式在对偏置线性化时关键（\cref{subsec:preint-bias}）。
\end{derivation}

\paragraph{连接恒等式与重力合并。}
$\mathbb{T}$ 与 $\mathrm{SE}_2(3)$ 的桥是
\begin{equation}
  \big(\boldsymbol{\mathcal{D}}(\boldsymbol{\Delta})\,\boldsymbol{\xi}\big)^\wedge=\boldsymbol{\Delta}^{-1}\boldsymbol{\xi}^\wedge\boldsymbol{\Delta},\qquad
  \mathrm{Exp}\big(\boldsymbol{\mathcal{D}}(\boldsymbol{\Delta})\boldsymbol{\xi}\big)=\boldsymbol{\Delta}^{-1}\mathrm{Exp}(\boldsymbol{\xi})\boldsymbol{\Delta},
  \label{eq:time-machine-conn}
\end{equation}
对级数逐项用前式即得后式。把 \cref{eq:preint-meanprop} 跨多步展开时，$\mathrm{SE}_2(3)$ 与 $\mathbb{T}$ 的元素交替出现；用 \cref{eq:time-machine-conn} 可自由移位，把\emph{所有}重力元素合并成一个 $\mathbf{T}_{g,J:0}$、所有时间机器合并成 $\boldsymbol{\Delta}_{J:0}=\boldsymbol{\Delta}(\tau_J-\tau_0)$\cite{barfoot2024state}：
\begin{equation}
  \mathbf{T}_{g,j:\ell}=\mathrm{Exp}\!\big(\big[\,\tfrac12(\tau_j-\tau_\ell)^2\mathbf{g}\;;\;(\tau_j-\tau_\ell)\mathbf{g}\;;\;\mathbf{0}\,\big]\big),
  \label{eq:gravity-merge}
\end{equation}
从而把 $J$ 步均值传播压成\emph{常数个}矩阵积、降低预积分代价。这正是下节误差定义 \cref{eq:preint-error} 里 $\mathbf{T}_{g,J:0}\boldsymbol{\Delta}_{J:0}^{-1}$ 的来历。

\subsection{预积分误差与端点扰动}
\label{subsec:preint-error}

\paragraph{误差定义。}
只在窗口两端 $\tau_0,\tau_J$ 保留状态 $\{\mathbf{T}_0,\mathbf{b}_0\},\{\mathbf{T}_J,\mathbf{b}_J\}$。预积分残差取“窗末状态”与“从窗首用 IMU 推得的预测”之差——扩展位姿部分用李群意义（取 $\mathrm{Log}$），偏置部分用加性\cite{barfoot2024state}：
\begin{equation}
  \mathbf{e}_T=\mathrm{Log}\!\Big(\mathbf{T}_J^{-1}\,\mathbf{T}_{g,J:0}\,\boldsymbol{\Delta}_{J:0}^{-1}\,\mathbf{T}_0\textstyle\prod_{j=1}^{J}\boldsymbol{\Delta}_j\boldsymbol{\Xi}_j\Big),\qquad
  \mathbf{e}_b=\begin{bmatrix}\mathbf{b}_{\omega,0}-\mathbf{b}_{\omega,J}\\\mathbf{b}_{a,0}-\mathbf{b}_{a,J}\end{bmatrix}.
  \label{eq:preint-error}
\end{equation}
$\mathbf{e}_T\in\mathbb{R}^9$、$\mathbf{e}_b\in\mathbb{R}^6$，合成 $15$ 维 $\mathbf{e}=[\mathbf{e}_T;\mathbf{e}_b]$。注意 $\boldsymbol{\Xi}_j$ \emph{只}依赖偏置与 IMU 读数、不依赖扩展位姿——这使端点扰动与偏置扰动可\emph{分别}线性化。

\begin{derivation}[端点扩展位姿扰动 $\to$ $\mathbf{V}_0$]
\label{deriv:preint-V0}
作右扰动 $\mathbf{T}_j=\mathbf{T}_{\mathrm{op},j}\mathrm{Exp}(\boldsymbol{\epsilon}_j)$（$j\in\{0,J\}$，与本节右扰动一致）代入 \cref{eq:preint-error}：
\[
  \mathbf{e}_T=\mathrm{Log}\!\Big(\mathrm{Exp}(-\boldsymbol{\epsilon}_J)\,\mathbf{T}_{\mathrm{op},J}^{-1}\mathbf{T}_{g,J:0}\boldsymbol{\Delta}_{J:0}^{-1}\mathbf{T}_{\mathrm{op},0}\,\mathrm{Exp}(\boldsymbol{\epsilon}_0)\textstyle\prod_j\boldsymbol{\Delta}_j\boldsymbol{\Xi}_j\Big).
\]
用伴随恒等式 $\mathbf{T}\,\mathrm{Exp}(\boldsymbol{\epsilon})=\mathrm{Exp}(\mathrm{Ad}(\mathbf{T})\boldsymbol{\epsilon})\,\mathbf{T}$（\cref{eq:se23-adjoint}）与时间机器连接 \cref{eq:time-machine-conn}，把 $\boldsymbol{\epsilon}_0$ 一路搬到最左：
\begin{equation}
  \mathbf{e}_T\approx\mathrm{Log}\!\Big(\mathrm{Exp}\!\big(\mathbf{V}_0\boldsymbol{\epsilon}_0-\boldsymbol{\epsilon}_J\big)\,\underbrace{\mathbf{T}_{\mathrm{op},J}^{-1}\mathbf{T}_{g,J:0}\boldsymbol{\Delta}_{J:0}^{-1}\mathbf{T}_{\mathrm{op},0}\textstyle\prod_j\boldsymbol{\Delta}_j\boldsymbol{\Xi}_j}_{\;=\;\mathrm{Exp}(\mathbf{e}_{\mathrm{op},T})}\Big),
  \label{eq:preint-eT-pert}
\end{equation}
其中（$\boldsymbol{\mathcal{T}}:=\mathrm{Ad}(\mathbf{T})$、$\boldsymbol{\mathcal{D}}_{J:0}:=\boldsymbol{\mathcal{D}}(\boldsymbol{\Delta}_{J:0})$、$\boldsymbol{\mathcal{T}}_{g,J:0}:=\mathrm{Ad}(\mathbf{T}_{g,J:0})$ 为常数）
\begin{equation}
  \mathbf{V}_0=\boldsymbol{\mathcal{T}}_{\mathrm{op},J}^{-1}\,\boldsymbol{\mathcal{T}}_{g,J:0}\,\boldsymbol{\mathcal{D}}_{J:0}\,\boldsymbol{\mathcal{T}}_{\mathrm{op},0}.
  \label{eq:preint-V0}
\end{equation}
端点扩展位姿的扰动\emph{只}经 $\mathbf{V}_0$（对窗首）与 $-\mathbf{I}$（对窗末）进入误差——简洁如线性系统。$\square$
\end{derivation}

\subsection{偏置扰动：\texorpdfstring{$\mathbf{J}^{-1}\mathbf{N}$}{J-inv N} 一阶展开}
\label{subsec:preint-bias}

偏置只藏在 $\boldsymbol{\Xi}_j$ 里（\cref{eq:preint-psi}），故对偏置线性化要展开 $\boldsymbol{\psi}_j$。难点在第一行的 $\mathbf{J}_j^{-1}\mathbf{N}_j$——$\mathbf{J}_j,\mathbf{N}_j$ 都依赖 $\bar{\boldsymbol{\omega}}_j=\tilde{\boldsymbol{\omega}}_j-\mathbf{b}_\omega$、从而依赖偏置。

\begin{derivation}[偏置雅可比 $\mathbf{B}_0$]
\label{deriv:preint-B0}
\textbf{Step 1（$\mathbf{J}^{-1}\mathbf{N}$ 的级数）.}\; 用 $\mathbf{J}(\boldsymbol{\phi})^{-1}=\sum_n\frac{B_n}{n!}(\boldsymbol{\phi}^\wedge)^n$（$B_n$ Bernoulli 数）与 \cref{eq:N-series}，两级数相乘得纯奇次幂闭式\cite{barfoot2024state}：
\begin{equation}
  \mathbf{J}(\boldsymbol{\phi})^{-1}\mathbf{N}(\boldsymbol{\phi})
  =\mathbf{I}-\tfrac16\boldsymbol{\phi}^\wedge+\tfrac{1}{360}(\boldsymbol{\phi}^\wedge)^3-\tfrac{1}{15120}(\boldsymbol{\phi}^\wedge)^5+\cdots
  \label{eq:JinvN}
\end{equation}
\textbf{Step 2（对偏置一阶展开）.}\; 设偏置扰动 $\mathbf{b}_{\omega,0}=\mathbf{b}_{\mathrm{op},\omega,0}+\boldsymbol{\beta}_{\omega,0}$、$\mathbf{b}_{a,0}=\mathbf{b}_{\mathrm{op},a,0}+\boldsymbol{\beta}_{a,0}$。代入 $\boldsymbol{\phi}=\Delta\tau_j\bar{\boldsymbol{\omega}}_{\mathrm{op},j}$（$\bar{\boldsymbol{\omega}}_{\mathrm{op},j}=\tilde{\boldsymbol{\omega}}_j-\mathbf{b}_{\mathrm{op},\omega,0}$），保留对 $\boldsymbol{\beta}$ 一阶、对 $\Delta\tau_j\bar{\boldsymbol{\omega}}$ 三阶：
\begin{equation}
  \mathbf{J}_j^{-1}\mathbf{N}_j\bar{\mathbf{a}}_j\approx\mathbf{J}_{\mathrm{op},j}^{-1}\mathbf{N}_{\mathrm{op},j}\bar{\mathbf{a}}_{\mathrm{op},j}
  -\tfrac16\Delta\tau_j\bar{\mathbf{a}}_{\mathrm{op},j}^\wedge\boldsymbol{\beta}_{\omega,0}
  +\tfrac{1}{360}\Delta\tau_j^3\mathbf{W}_j\boldsymbol{\beta}_{\omega,0}
  -\mathbf{J}_{\mathrm{op},j}^{-1}\mathbf{N}_{\mathrm{op},j}\boldsymbol{\beta}_{a,0},
  \label{eq:preint-bias-firstorder}
\end{equation}
其中 $\mathbf{W}_j=\bar{\boldsymbol{\omega}}_{\mathrm{op},j}^{\wedge2}\bar{\mathbf{a}}_{\mathrm{op},j}^\wedge+\bar{\boldsymbol{\omega}}_{\mathrm{op},j}^\wedge(\bar{\boldsymbol{\omega}}_{\mathrm{op},j}^\wedge\bar{\mathbf{a}}_{\mathrm{op},j})^\wedge+(\bar{\boldsymbol{\omega}}_{\mathrm{op},j}^{\wedge2}\bar{\mathbf{a}}_{\mathrm{op},j})^\wedge$。
\textbf{Step 3（装成 $\boldsymbol{\psi}_j$ 的偏置块 $\mathbf{B}_j$）.}\; 于是 $\boldsymbol{\psi}_j\approx\boldsymbol{\psi}_{\mathrm{op},j}+\mathbf{B}_j\boldsymbol{\beta}_0$（$\boldsymbol{\beta}_0=[\boldsymbol{\beta}_{\omega,0};\boldsymbol{\beta}_{a,0}]$），按 $[\,\text{位置};\text{速度};\text{旋转}\,]$ 行序、$[\boldsymbol{\beta}_\omega;\boldsymbol{\beta}_a]$ 列序：
\begin{equation}
  \mathbf{B}_j=\begin{bmatrix}
  -\tfrac{1}{12}\Delta\tau_j^3\bar{\mathbf{a}}_{\mathrm{op},j}^\wedge+\tfrac{1}{720}\Delta\tau_j^5\mathbf{W}_j & -\tfrac12\Delta\tau_j^2\mathbf{J}_{\mathrm{op},j}^{-1}\mathbf{N}_{\mathrm{op},j}\\
  \mathbf{0} & -\Delta\tau_j\mathbf{I}\\
  -\Delta\tau_j\mathbf{I} & \mathbf{0}
  \end{bmatrix}.
  \label{eq:preint-Bj}
\end{equation}
（为省算可略 $\Delta\tau_j^3,\Delta\tau_j^5$ 项。）\textbf{Step 4（沿窗口累积 $\to$ $\mathbf{B}_0$）.}\; 把单步偏置扰动 $\mathrm{Exp}(\boldsymbol{\psi}_{\mathrm{op},j})\mathrm{Exp}(\boldsymbol{\mathcal{J}}(-\boldsymbol{\psi}_{\mathrm{op},j})\mathbf{B}_j\boldsymbol{\beta}_0)$（$\boldsymbol{\mathcal{J}}$ 为 $\mathrm{SE}_2(3)$ 左雅可比）沿 $\prod_j\boldsymbol{\Delta}_j\boldsymbol{\Xi}_j$ 全部搬到最左，用 \cref{eq:time-machine-conn} 与伴随律得 $\prod_j\boldsymbol{\Delta}_j\boldsymbol{\Xi}_j\approx\mathrm{Exp}(\mathbf{B}_0\boldsymbol{\beta}_0)\prod_j\boldsymbol{\Delta}_j\boldsymbol{\Xi}_{\mathrm{op},j}$，其中
\begin{equation}
  \mathbf{B}_0=\sum_{j=1}^{J}\Big(\prod_{\ell=1}^{j}\boldsymbol{\mathcal{D}}_\ell^{-1}\,\mathrm{Ad}(\boldsymbol{\Xi}_{\mathrm{op},\ell})\Big)\boldsymbol{\mathcal{J}}(-\boldsymbol{\psi}_{\mathrm{op},j})\,\mathbf{B}_j.
  \label{eq:preint-B0}
\end{equation}
$\square$
\end{derivation}

\subsection{装配线性化误差与全运动先验代价}
\label{subsec:preint-assemble}

\paragraph{$15$ 维端点雅可比。}
把 \cref{eq:preint-B0} 的偏置项与 \cref{deriv:preint-V0} 的端点项一并代回 \cref{eq:preint-eT-pert}，对 $\mathrm{Log}$ 用 $\mathrm{SE}_2(3)$ 左雅可比逆 $\boldsymbol{\mathcal{J}}(\mathbf{e}_{\mathrm{op},T})^{-1}$ 一阶展开，得线性化扩展位姿误差
\begin{equation}
  \mathbf{e}_T\approx\mathbf{e}_{\mathrm{op},T}+\boldsymbol{\mathcal{J}}(\mathbf{e}_{\mathrm{op},T})^{-1}\big(\mathbf{V}_0\boldsymbol{\epsilon}_0-\boldsymbol{\epsilon}_J+\mathbf{V}_0\mathbf{B}_0\boldsymbol{\beta}_0\big).
  \label{eq:preint-eT-lin}
\end{equation}
合上偏置误差 $\mathbf{e}_b$，整 $15$ 维残差对窗首/窗末扰动 $\boldsymbol{\varepsilon}_0=[\boldsymbol{\epsilon}_0;\boldsymbol{\beta}_0]$、$\boldsymbol{\varepsilon}_J=[\boldsymbol{\epsilon}_J;\boldsymbol{\beta}_J]$ 的雅可比为
\begin{equation}
  \mathbf{e}=\mathbf{e}_{\mathrm{op}}+\mathbf{F}_0\,\boldsymbol{\varepsilon}_0-\mathbf{E}_J\,\boldsymbol{\varepsilon}_J,\quad
  \mathbf{F}_0=\begin{bmatrix}\boldsymbol{\mathcal{J}}^{-1}\mathbf{V}_0&\boldsymbol{\mathcal{J}}^{-1}\mathbf{V}_0\mathbf{B}_0\\\mathbf{0}&\mathbf{I}\end{bmatrix},\;
  \mathbf{E}_J=\begin{bmatrix}\boldsymbol{\mathcal{J}}^{-1}&\mathbf{0}\\\mathbf{0}&\mathbf{I}\end{bmatrix},
  \label{eq:preint-FE}
\end{equation}
简记 $\boldsymbol{\mathcal{J}}^{-1}:=\boldsymbol{\mathcal{J}}(\mathbf{e}_{\mathrm{op},T})^{-1}$。这与 \cref{ch:nlopt} 的运动先验因子 $\mathbf{e}_{v,k}=\mathbf{e}_{v,k}(\mathbf{x}_{\mathrm{op}})+\mathbf{F}_{k-1}\boldsymbol{\varepsilon}_{k-1}-\mathbf{E}_k\boldsymbol{\varepsilon}_k$ \emph{同型}——预积分把一窗 IMU 变成了一条标准“双端”因子。

\paragraph{拼成因子图代价。}
把局部端点 $\tau_0,\tau_J$ 换成全局时刻 $t_{k-1},t_k$，叠加窗首先验 $\{\check{\mathbf{T}}_0,\check{\mathbf{b}}_0\}$（协方差 $\check{\mathbf{P}}_0$），按 \cref{ch:nlopt} 堆叠各窗误差 $\mathbf{e}_v(\mathbf{x}_{\mathrm{op}})$、块对角权重 $\mathbf{Q}=\mathrm{diag}(\check{\mathbf{P}}_0,\mathbf{Q}_1,\dots,\mathbf{Q}_K)$ 与下双对角“逆转移”
\begin{equation}
  \mathbf{F}^{-1}=\begin{bmatrix}\mathbf{E}_0&&&\\-\mathbf{F}_0&\mathbf{E}_1&&\\&\ddots&\ddots&\\&&-\mathbf{F}_{K-1}&\mathbf{E}_K\end{bmatrix},
  \label{eq:preint-Finv}
\end{equation}
则运动先验代价为
\begin{equation}
  J_v\approx\tfrac12\big(\mathbf{e}_v(\mathbf{x}_{\mathrm{op}})-\mathbf{F}^{-1}\delta\mathbf{x}\big)^\top\mathbf{Q}^{-1}\big(\mathbf{e}_v(\mathbf{x}_{\mathrm{op}})-\mathbf{F}^{-1}\delta\mathbf{x}\big),
  \label{eq:preint-cost}
\end{equation}
可与任意观测代价相加成完整 MAP 问题。每窗协方差 $\mathbf{Q}_k$ 由 \cref{deriv:Q-series} 的离散过程噪声在窗内逐小步复合而成（\cref{eq:Q-batch} 正是其“预测误差二阶矩”依据）；$\mathbf{Q}_k$ 依赖该窗偏置与 IMU 读数、\emph{不}依赖扩展位姿，故每轮迭代按上轮偏置重算。每轮更新用各自扰动方案 $\mathbf{T}_{\mathrm{op},k}\leftarrow\mathbf{T}_{\mathrm{op},k}\mathrm{Exp}(\boldsymbol{\epsilon}_k^\star)$、$\mathbf{b}_{\mathrm{op},k}\leftarrow\mathbf{b}_{\mathrm{op},k}+\boldsymbol{\beta}_k^\star$，迭代至收敛。

\subsection{左右扰动互换与扩展位姿 \texorpdfstring{$\odot$}{odot} 算子}
\label{subsec:preint-leftright}

把上述预积分先验接到\emph{观测}（如点云/路标）时，需扩展位姿版的 $\odot$ 算子。本书 \cref{ch:lie} 仅给 $\SE(3)$ 的 $4\times6$ $\odot$；扩展位姿要把\emph{速度块置零}以“正确忽略速度分量”\cite{barfoot2024state}。对齐次路标 $\mathbf{p}=[\mathbf{q};1;0]$（最后一位 $0$ 屏蔽速度槽），
\begin{equation}
  \mathbf{p}^\odot=\begin{bmatrix}\mathbf{q}\\1\\0\end{bmatrix}^{\!\odot}=\begin{bmatrix}\mathbf{I}&-\mathbf{q}^\wedge&\mathbf{0}\\\mathbf{0}^\top&\mathbf{0}^\top&\mathbf{0}^\top\\\mathbf{0}^\top&\mathbf{0}^\top&\mathbf{0}^\top\end{bmatrix}\in\mathbb{R}^{5\times9},
  \label{eq:se23-odot}
\end{equation}
满足 $\boldsymbol{\epsilon}^\wedge\mathbf{p}=\mathbf{p}^\odot\boldsymbol{\epsilon}$（列序 $[\,\text{位置};\text{速度};\text{旋转}\,]$）。中间 $-\mathbf{q}^\wedge$ 在\emph{旋转}列、位置列为 $\mathbf{I}$、\textbf{速度列恒 $\mathbf{0}$}——这正是 \cref{eq:inekf-H} 观测雅可比里“速度块为零”的来历。于是对观测 $\mathbf{y}=\mathbf{D}^\top\mathbf{T}_{iv}^{-1}\mathbf{p}$ 作右扰动 $\mathbf{T}_{iv}=\mathbf{T}_{\mathrm{op},iv}\mathrm{Exp}(\boldsymbol{\epsilon}_r)$ 得雅可比 $\mathbf{G}_r=-\mathbf{D}^\top(\mathbf{T}_{\mathrm{op},iv}^{-1}\mathbf{p})^\odot$；左扰动得 $\mathbf{G}_\ell=+\mathbf{D}^\top(\mathbf{T}_{\mathrm{op},vi}\mathbf{p})^\odot$。因 $\boldsymbol{\epsilon}_r=-\boldsymbol{\epsilon}_\ell$，故 $\mathbf{G}_r=-\mathbf{G}_\ell$——\emph{算出一个即得另一个}（与 \cref{subsec:c4-bridge} 的左右桥同源）。

\begin{insight}[$\mathrm{SE}_2(3)$ 预积分 vs Forster 预积分]
两者都把一窗 IMU 压成单因子，区别在\emph{误差坐标}：Forster\cite{forster2017preintegration} 在 $\SO(3)\times\mathbb{R}^3\times\mathbb{R}^3$ 上分量式定义 $\Delta\mathbf{R},\Delta\mathbf{v},\Delta\mathbf{p}$，偏置修正用三块独立一阶雅可比；$\mathrm{SE}_2(3)$ 版（本节）把三者并进\emph{一个}群元素，端点雅可比 $\mathbf{F}_0,\mathbf{E}_J$（\cref{eq:preint-FE}）经\emph{一个} $\boldsymbol{\mathcal{J}}^{-1}$ 统一搬运、偏置经\emph{一个} $\mathbf{B}_0$（\cref{eq:preint-B0}）。代价是要 $\mathbf{N}$ 矩阵与时间机器等额外机械，收益是与本章 InEKF \emph{共享}全部 $\mathrm{SE}_2(3)$ 零件、且误差几何更一致（群仿射的批量映像）。工程选型：高频 IMU $+$ 成熟库 $\to$ Forster；要与 InEKF 滤波器同栈、或研究一致性 $\to$ $\mathrm{SE}_2(3)$ 版。
\end{insight}

\begin{practice}
\begin{enumerate}
  \item 由 \cref{eq:N-series} 验证 $\mathbf{N}(\mathbf{0})=\mathbf{I}$，并对小 $\phi$ 展开 \cref{eq:N-closed} 到 $\mathcal{O}(\phi^2)$，确认与级数前两项一致。
  \item 验证时间机器共轭 $\boldsymbol{\Delta}_j^{-1}\mathbf{T}_{j-1}\boldsymbol{\Delta}_j$（\cref{deriv:preint-triple}）确为 $\big[\begin{smallmatrix}\mathbf{R}&\mathbf{r}+\Delta\tau\mathbf{v}&\mathbf{v}\\\mathbf{0}^\top&1&0\\\mathbf{0}^\top&0&1\end{smallmatrix}\big]$，并说明它如何吸收 \cref{eq:preint-Rprop} 的跨步耦合 $\Delta\tau_j\bar{\mathbf{v}}(\tau_{j-1})$。
  \item 取 $J=1$（单步预积分），写出 \cref{eq:preint-error} 退化形式，并对照 \cref{eq:se23-Phi} 的幂零转移矩阵，指出 $\mathbf{V}_0$（\cref{eq:preint-V0}）与 $\boldsymbol{\Phi}(\Delta t)$ 的关系。
  \item （概念）解释为何 $\mathbf{Q}_k$ 不依赖扩展位姿、却依赖偏置——这对“每轮重算 $\mathbf{Q}_k$”的实现意味着什么？
\end{enumerate}
\end{practice}

% ════════════════════════════════════════════════════════════════
\section*{本章常见误解汇总}
\addcontentsline{toc}{section}{本章常见误解汇总}

\begin{table}[H]\centering\small
\begin{tabular}{p{0.42\linewidth}p{0.52\linewidth}}
\toprule
误解 & 正确理解 \\
\midrule
四元数归一化后就能直接用标准 EKF & 归一化只修模值；$4\times4$ 协方差沿径向仍零特征值、Cholesky 必奇异。正解用 $3$ 维误差协方差（MEKF/ESKF），见 \cref{sec:why-manifold} \\
InEKF 只是“把 EKF 放到李群上” & 流形化 ESKF 已做到；InEKF 的核心是\emph{群仿射 $+$ 不变误差}使 $\mathbf{A}$ 与估计无关（\cref{thm:log-linear}），从根消除虚假可观测性——结构性优势而非记号变换 \\
log-linear 是线性化近似 & \cref{thm:log-linear} 断言对数坐标误差\emph{精确}满足 $\dot{\boldsymbol{\xi}}=\mathbf{A}_t\boldsymbol{\xi}$（无噪声传播段），BCH 交叉项经群同态 $+\,[\mathbf{X},\mathbf{X}]=\mathbf{0}$ 完全消除（\cref{deriv:log-linear}） \\
InEKF 完全消除所有非线性 & 仅\emph{传播段}精确；更新步仍有二阶余项 $\mathbf{r}_n=\mathcal{O}(\|\boldsymbol{\xi}\|\|\mathbf{L}\boldsymbol{\xi}\|)$。\cref{thm:inekf-stable} 的非平凡处正是论证该余项一致有界 \\
group-affine 就是左不变加右不变 & group-affine 严格\emph{扩张}左右不变的并集；判据看 $f(\mathbf{e})$ 是否为零，$f(\mathbf{e})\neq\mathbf{0}$（重力）才是新意（\cref{sec:group-affine}） \\
零偏可用 group-affine 处理 & negative result（\cref{thm:bias-negative}）：$\boldsymbol{\omega}-\mathbf{b}_g$ 引入 $\mathbf{b}_1^\wedge\neq(\mathbf{R}_2^\top\mathbf{b}_1)^\wedge$ 矛盾。须用不完美 InEKF/切群/EqF \\
左/右不变误差只是符号惯例 & 由\emph{观测}的不变性决定（\cref{subsec:obs-classify}）：landmark 右不变、GPS 左不变。手性错配使 $\mathbf{H}$ 重新依赖估计、\cref{thm:inekf-stable} 失效 \\
InEKF 自动全局收敛 & \cref{thm:inekf-stable} 仍是\emph{局部}收敛，只是 $\varepsilon\perp t_0$；受 $\log$ 单射半径限制 \\
IKFoM（FAST-LIO2）是 InEKF & IKFoM 是 \emph{iterated} EKF on manifold；状态是 product manifold、$\mathbf{F}$ 依赖 $\hat{\mathbf{R}}$、不享 log-linear（\cref{subsec:ikfom}） \\
UKF-M 免雅可比所以更好 & 免的是\emph{人力}（不手推），计算成本更高且不享 InEKF 的“误差动力学与状态解耦”；两者互补 \\
\bottomrule
\end{tabular}
\end{table}

% ════════════════════════════════════════════════════════════════
\section*{本章小结}
\addcontentsline{toc}{section}{本章小结}

\paragraph{符号表。}

\begin{table}[H]\centering\small
\caption{本章主要符号}
\begin{tabular}{lll}
\toprule
符号 & 含义 & 出处 \\
\midrule
$\mathbf{X}\in G$ & 李群状态（源文献记 $\chi$） & \cref{sec:group-affine} \\
$\hat{\mathbf{X}},\mathbf{X}$ & 估计 / 真值轨迹 & \cref{eq:lr-error} \\
$\boldsymbol{\eta}^L=\mathbf{X}^{-1}\hat{\mathbf{X}}$ & 左不变误差（估计在右，TAC 2017 约定） & \cref{eq:lr-error} \\
$\boldsymbol{\eta}^R=\hat{\mathbf{X}}\mathbf{X}^{-1}$ & 右不变误差（同 \cref{paper:inekf}） & \cref{eq:lr-error} \\
$\boldsymbol{\xi}=\mathrm{Log}(\boldsymbol{\eta})$ & 误差的对数坐标 & \cref{thm:log-linear} \\
$f_{\mathbf{u}}(\mathbf{X})$ & 群上动力学（输入 $\mathbf{u}$） & \cref{eq:ga-dynamics} \\
$\mathbf{A}_{\mathbf{u}_t}$ & log-linear 系统矩阵（仅依赖输入） & \cref{eq:log-linear} \\
$\boldsymbol{\Phi}$ & 误差状态转移（$\mathrm{SE}_2(3)$ 幂零闭式） & \cref{eq:se23-Phi} \\
$\mathbf{H}$ & 不变观测雅可比（不依赖估计） & \cref{eq:inekf-H} \\
$\mathbf{P},\mathbf{K},\mathbf{S}$ & 协方差 / 增益 / 新息协方差 & \cref{algo:riekf} \\
$\delta\boldsymbol{\theta},\delta\mathbf{v},\delta\mathbf{p}$ & 姿态 / 速度 / 位置误差 & \cref{deriv:mekf-theta} \\
$\mathbf{b}_g,\mathbf{b}_a$ & 陀螺 / 加速度计零偏 & \cref{sec:bias-negative} \\
\bottomrule
\end{tabular}
\end{table}

\paragraph{定理速查表。}

\begin{table}[H]\centering\small
\caption{本章定理速查}
\begin{tabular}{lll}
\toprule
名称 & 核心结论 & 标签 \\
\midrule
MEKF/ESKF 角误差 & $\dot{\delta\boldsymbol{\theta}}=-\hat{\boldsymbol{\omega}}^\wedge\delta\boldsymbol{\theta}-\delta\mathbf{b}_g-\mathbf{n}_g$ & \cref{deriv:mekf-theta} \\
群仿射三等价 & 左/右不变误差自治 $\Leftrightarrow$ $f(\mathbf{a}\mathbf{b})=f(\mathbf{a})\mathbf{b}+\mathbf{a}f(\mathbf{b})-\mathbf{a}f(\mathbf{e})\mathbf{b}$ & \cref{thm:ga-equiv} \\
$\mathrm{SE}_2(3)$ IMU 群仿射 & $f(\mathbf{X})=\mathbf{A}\mathbf{X}+\mathbf{X}\mathbf{B}$，逐块验证 & \cref{deriv:se23-ga} \\
log-linear & $\dot{\boldsymbol{\xi}}=\mathbf{A}_{\mathbf{u}_t}\boldsymbol{\xi}$ 精确，$\mathbf{A}_{\mathbf{u}_t}\perp\hat{\mathbf{X}}$ & \cref{thm:log-linear} \\
$\mathbf{A}$ 幂零 & $\mathbf{A}^3=\mathbf{0}$，$\boldsymbol{\Phi}=\mathbf{I}+\mathbf{A}\Delta t+\tfrac12\mathbf{A}^2\Delta t^2$ 精确 & \cref{eq:se23-A,eq:se23-Phi} \\
零偏 negative result & 不存在李群使 pose$+$bias 群仿射；$\mathbf{b}_1^\wedge\neq(\mathbf{R}_2^\top\mathbf{b}_1)^\wedge$ & \cref{thm:bias-negative} \\
InEKF 稳定性 & 渐近稳定，收敛半径 $\varepsilon\perp t_0$ & \cref{thm:inekf-stable} \\
$\mathbf{N}$ 矩阵 & $\mathbf{N}(\phi\mathbf{a})=2\!\int_0^1\!\alpha\mathbf{J}(\alpha\phi\mathbf{a})d\alpha$，位置均值精确积分 & \cref{def:N-matrix} \\
$\mathrm{SE}_2(3)$ 预积分 & $\mathbf{e}_T=\mathrm{Log}(\cdots)$，端点雅可比 $\mathbf{F}_0/\mathbf{E}_J$ 仅经一个 $\boldsymbol{\mathcal{J}}^{-1}$ & \cref{eq:preint-error,eq:preint-FE} \\
\bottomrule
\end{tabular}
\end{table}

\paragraph{知识点总表（难度 $\bigstar$）。}

\begin{table}[H]\centering\small
\caption{本章知识点与难度}
\begin{tabular}{llll}
\toprule
\# & 知识点 & 难度 & 地位 \\
\midrule
1 & 三类结构性失败（奇异/奇点/伪观测） & $\bigstar\bigstar$ & 动机基石 \\
2 & 统一原则：名义$+$误差、$\boxplus/\boxminus$ & $\bigstar\bigstar$ & 主干 \\
3 & MEKF $\to$ ESKF 乘性误差推导 & $\bigstar\bigstar\bigstar$ & 承接 \cref{ch:eskf} \\
4 & 左/右不变误差与约定坑 & $\bigstar\bigstar\bigstar$ & InEKF 入口 \\
5 & 群仿射三等价定理 $+$ 证明 & $\bigstar\bigstar\bigstar\bigstar$ & 理论核心 \\
6 & $\mathrm{SE}_2(3)$ 逐块验证群仿射 & $\bigstar\bigstar\bigstar$ & 落地核心 \\
7 & log-linear 精确性 $+$ BCH 机理 & $\bigstar\bigstar\bigstar\bigstar$ & 理论巅峰 \\
8 & RIEKF 完整离散算法 & $\bigstar\bigstar\bigstar$ & 工程主线 \\
9 & 一致性：$\mathbf{H}\mathbf{N}=\mathbf{0}$ 结构保持 & $\bigstar\bigstar\bigstar\bigstar$ & 最大卖点 \\
10 & 零偏 negative result $+$ 修复路径 & $\bigstar\bigstar\bigstar$ & 边界 \\
11 & UKF-M / IKFoM / EqF 全族 & $\bigstar\bigstar\bigstar$ & 广度 \\
12 & 稳定性 $\varepsilon\perp t_0$ & $\bigstar\bigstar\bigstar\bigstar$ & 选读巅峰 \\
13 & $\mathrm{SE}_2(3)$ 不变预积分（$\mathbf{N}$/时间机器/$\mathbf{F}_0\mathbf{E}_J$） & $\bigstar\bigstar\bigstar\bigstar$ & 批量收尾 capstone \\
\bottomrule
\end{tabular}
\end{table}

% ════════════════════════════════════════════════════════════════
\section*{流形滤波族全景对比}
\addcontentsline{toc}{section}{流形滤波族全景对比}

\begin{table}[H]\centering\small
\caption{流形滤波族横评}
\label{tab:family-compare}
\begin{tabular}{lllll}
\toprule
方法 & 状态空间 & 误差定义 & 雅可比需求 & 一致性 \\
\midrule
标准 EKF & $\mathbb{R}^n$（含四元数） & 加性 & 解析，沿轨迹 & 过自信（FEJ 修补） \\
MEKF & $\SO(3)$ $+$ 欧氏 bias & 姿态乘性 $+$ bias 加性 & 解析（简单） & 良好（姿态） \\
ESKF & $\SE(3)$/IMU $+$ 欧氏 & 姿态乘性 $+$ 其余加性 & 解析（$\mathbf{F}$ 矩阵） & 良好（需复位） \\
InEKF（完美） & 李群 $G$（$\mathrm{SE}_2(3)$） & $\boldsymbol{\eta}=\hat{\mathbf{X}}\mathbf{X}^{-1}$ & \emph{精确、$\perp\hat{\mathbf{X}}$} & \emph{自动一致}（log-linear） \\
不完美 InEKF & $G\times\mathbb{R}^k$ & pose 乘性 $+$ bias 加性 & $\mathbf{A}$ 部分依赖 $\hat{\mathbf{X}}$ & 弱化、经验优越 \\
IKFoM & 通用混合流形 & $\boxplus/\boxminus$ & 解析（用户提供） & 局部 \\
UKF-M & 可平行化流形 & retraction $\varphi$ & \emph{免解析}（sigma） & 取决于 $\varphi$ \\
EqF & 齐性空间 $G/H$ & 固定原点等变误差 & 线性化于原点 & 严格一致 \\
\bottomrule
\end{tabular}
\end{table}

\paragraph{选型决策。}\emph{不可能三角}（类比 CAP）：一致性、通用性、实现简单性难以兼得。$\mathbf{(1)}$ 只有姿态 $\to$ MEKF；$\mathbf{(2)}$ IMU $+$ 外部观测、要工程主流 $\to$ ESKF；$\mathbf{(3)}$ 动力学群仿射、无长期视觉零偏（腿足、无 bias 惯导）$\to$ InEKF；$\mathbf{(4)}$ 雅可比难算 $\to$ UKF-M；$\mathbf{(5)}$ 通用混合流形 $+$ LiDAR $\to$ IKFoM（FAST-LIO2）；$\mathbf{(6)}$ 齐性空间 $+$ 要严格一致（单目 VIO）$\to$ EqF。\textbf{所有选择的前提是同一件事——把扰动约定、切空间排序、四元数代数像宪法一样贯穿整个代码库}（流形滤波集成 bug 的头号来源）。

\begin{note}
\textbf{扰动/排序约定速查（集成 bug 头号源）.}\;
manif 切空间 $[\boldsymbol{\rho},\boldsymbol{\theta}]$（trans-first）、Sophus $[\boldsymbol{\upsilon},\boldsymbol{\omega}]$（trans-first）、\textbf{GTSAM $[\boldsymbol{\omega},\boldsymbol{\rho}]$（rot-first，易错！）}、OpenVINS 用 JPL 四元数 $+$ 左乘误差、本书与 Barfoot $\mathrm{SE}_2(3)$ 用 $[\boldsymbol{\rho};\boldsymbol{\nu};\boldsymbol{\phi}]$（\cref{def:se23}）而 Barrau/Hartley 用 $[\boldsymbol{\phi};\boldsymbol{\nu};\boldsymbol{\rho}]$。两大陷阱：(1) GTSAM\,$\leftrightarrow$\,Sophus/manif 的 $6$D 顺序互反，协方差直接互喂会让位置/姿态椭球\emph{对调}，需置换矩阵 $\boldsymbol{\Pi}=\big[\begin{smallmatrix}\mathbf{0}&\mathbf{I}_3\\\mathbf{I}_3&\mathbf{0}\end{smallmatrix}\big]$ 重排 $\mathbf{P}'=\boldsymbol{\Pi}\mathbf{P}\boldsymbol{\Pi}^\top$；(2) Hamilton（$ij=k$，Eigen/Sophus/GTSAM/ROS/VINS-Mono）$\leftrightarrow$ JPL（$ji=k$，OpenVINS/MSCKF/AOCS）相同 $(x,y,z,w)$ 数值语义\emph{相反}，须显式转换层。
\end{note}

% ════════════════════════════════════════════════════════════════
\section*{延伸阅读}
\addcontentsline{toc}{section}{延伸阅读}

\begin{itemize}
  \item \textbf{奠基论文}：Barrau \& Bonnabel《The Invariant Extended Kalman Filter as a Stable Observer》\cite{barrau2017invariant}——群仿射、log-linear、稳定性的源头（精读 §II$+$附录 B/C）；Bonnabel--Martin--Rouchon《Symmetry-Preserving Observers》\cite{bonnabel2008symmetry}——不变观测器的历史框架（Cartan moving frame）。
  \item \textbf{ESKF/MEKF}：Solà《Quaternion kinematics for the error-state Kalman filter》\cite{sola2017quaternion}（$73$ 页 ESKF 完整推导）；Solà《A micro Lie theory》\cite{sola2018micro}（$\boxplus/\boxminus$ 统一）；Lefferts--Markley--Shuster\cite{lefferts1982kalman}、Markley\cite{markley2003attitude}（MEKF 原型与误差表示对比）。
  \item \textbf{工程落地}：Hartley 等《Contact-Aided Invariant EKF》\cite{hartley2020contact}（Cassie 接触辅助 InEKF $+$ 不完美 InEKF）；Brossard 等 UKF-M\cite{brossard2020ukfm}；He--Xu--Zhang IKFoM\cite{he2021ikfom} 与 Xu 等 FAST-LIO2\cite{xu2022fastlio2}；Brossard 等 $\mathrm{SE}_2(3)$ 不变预积分\cite{brossard2021associating}。
  \item \textbf{前沿推广}：Mahony--Trumpf EqF\cite{mahony2021equivariant}、van Goor--Mahony EqVIO\cite{vangoor2023eqvio}、Fornasier 等零偏 EqF\cite{fornasier2022overcoming}；Deyst--Price 1968\cite{deyst1968conditions}（LTV-KF 稳定性原始结果）；Chirikjian Vol.2\cite{chirikjian2011stochastic}（BCH 背景）；Barfoot 第 2 版\cite{barfoot2024state} §8.4（工程实现）。
\end{itemize}

% ════════════════════════════════════════════════════════════════
\section*{与后续章节的关系}
\addcontentsline{toc}{section}{与后续章节的关系}

\begin{itemize}
  \item \textbf{$\leftarrow$ \cref{ch:lie}（李群）}：本章的 $\mathrm{SE}_2(3)$（\cref{def:se23,thm:se23-exp,eq:se23-adjoint}）、对称/不变/等变（\cref{sec:lie-symmetry,def:invariant-error}）、$\exp/\log/\mathrm{Ad}$ 全来自李群章，是 InEKF 的代数地基。
  \item \textbf{$\leftarrow$ \cref{ch:eskf}（ESKF）}：本章是其 \cref{paper:inekf} 的完整展开；ESKF（\cref{sec:eskf}）、KF（\cref{thm:kf}）、一致性诊断（\cref{def:nees-nis}）、滤波器谱系（\cref{tab:filter-taxo}）是直接前置。
  \item \textbf{$\to$ \cref{ch:imu}（IMU）}：$\mathrm{SE}_2(3)$ 不变预积分（\cref{sec:preint-se23} 已给完整推导）\cite{brossard2021associating} 把本章群仿射结构用于协方差精确传播，是 Forster 预积分\cite{forster2017preintegration} 的范式升级；\cref{ch:imu} 给主流 Forster $\SO(3)$ 版作对照。
  \item \textbf{$\to$ \cref{ch:vio}（VIO）}：InEKF/EqF 是当代 VIO 一致性研究的主流框架；本章一致性分析（\cref{sec:consistency-inekf}）解释“右不变卡尔曼滤波”去向。
  \item \textbf{$\to$ 腿足机器人状态估计}：接触辅助 InEKF（\cref{sec:bias-negative}）直接服务于运动控制的状态反馈（接触辅助预积分\cite{hartley2018legged}、全地形腿足里程计\cite{wisth2023vilens}）。
\end{itemize}

% ════════════════════════════════════════════════════════════════
\section*{故障排查手册}
\addcontentsline{toc}{section}{故障排查手册}

\begin{table}[H]\centering\small
\begin{tabularx}{\linewidth}{LLL}
\toprule
症状 & 可能原因 & 排查步骤 \\
\midrule
量测更新后状态\emph{远离}真值 & residual 符号或左右乘方向错 & 在 $\mathbb{R}$ 上构造最小反例；核对“误差定义—residual—乘向”三者闭合（\cref{algo:riekf} 第 3 步） \\
传播 Jacobian 仍依赖估计 & 系统不群仿射或观测手性选错 & 逐块验证群仿射残差（\cref{deriv:se23-ga}）；检查 $\mathbf{Y}=\mathbf{X}\mathbf{d}$ vs $\mathbf{X}^{-1}\mathbf{d}$（\cref{subsec:obs-classify}） \\
加入零偏后 NEES 发散 & 零偏破坏群仿射（\cref{thm:bias-negative}） & 改用不完美 InEKF；增大 $\mathbf{Q}_{\text{bias}}$；考虑 EqF \\
旋转误差接近 $\pi$ 滤波器跳变 & $\log$ 在 $\pi$ 处多值/不连续 & 检查四元数是否跨对径；降低初始化误差；对相对旋转取 $\mathrm{Log}$ \\
VIO yaw 方差异常收缩 & 线性化破坏不可观结构 & 计算 $\mathbf{H}\mathbf{N}$ 是否为零（\cref{eq:obs-null}）；对比 FEJ 开关；检查 InEKF 误差定义 \\
大步长传播 NEES 异常 & 状态/噪声积分近似过粗 & 更细 IMU 积分或 Van Loan 离散；勿误以为 $\exp(\mathbf{A}\Delta t)$ 本身破坏群仿射 \\
不同库输出姿态相反 & Hamilton/JPL 或左右扰动混用 & 同一 $\mathbf{q}$ 作用已知向量对比；检查 $\mathrm{Ad}$ 方向；加显式转换层 \\
Cholesky 失败 & 协方差失正定 & Joseph 形式（\cref{algo:riekf} 第 4 步）；对称化 $\mathbf{P}\leftarrow\tfrac12(\mathbf{P}+\mathbf{P}^\top)$ \\
\bottomrule
\end{tabularx}
\end{table}

\begin{insight}
\textbf{贯穿本章的一句话}：流形滤波的力量来自\emph{对称性}，而不是来自“在李群上写 EKF”。只有\emph{动力学、误差与观测三者共同尊重群结构}，log-linear 定理才真正成立；只要误差定义、创新量与注入方向三者不闭合，同一套增益公式就会把正确量测变成错误校正。这比任何高深理论都更决定你的 VIO/惯导/腿足系统能否活过真实场景的前十分钟。
\end{insight}
